\documentclass[11pt,a4paper]{article}
\usepackage[margin=23mm]{geometry}
\usepackage[T1]{fontenc}\usepackage[utf8]{inputenc}
\usepackage{lmodern,microtype,amsmath,amssymb,amsthm,mathtools,mathrsfs}
\usepackage{booktabs,longtable,tabularx,array,enumitem}
\usepackage[hidelinks]{hyperref}
\hypersetup{pdftitle={Exact degree and real-point counts in a localized sector of QLC connections},pdfauthor={Marius Torsheim}}\usepackage{fancyhdr}
\setlength{\headheight}{14pt}\setlength{\emergencystretch}{3em}
\pagestyle{fancy}\fancyhf{}\lhead{Exact enumeration on the binary cube}
\rhead{Quantum Levi--Civita connections}\cfoot{\thepage}
\newtheorem{theorem}{Theorem}[section]
\newtheorem{lemma}[theorem]{Lemma}\newtheorem{proposition}[theorem]{Proposition}
\newtheorem{corollary}[theorem]{Corollary}\newtheorem{definition}[theorem]{Definition}
\theoremstyle{remark}\newtheorem{remark}[theorem]{Remark}
\newcommand{\Spec}{\operatorname{Spec}}\newcommand{\tr}{\operatorname{tr}}
\newcommand{\Id}{\mathrm I}\newcommand{\F}{\mathbb F}
\newcommand{\Ab}{\overline A}\newcommand{\Bb}{\overline B}
\newcommand{\Orb}{\mathcal O}\newcommand{\length}{\operatorname{length}}
\newcommand{\norm}[1]{\left\|#1\right\|_\infty}
\title{\textbf{Exact degree and real-point counts in a localized sector of}\\
\textbf{quantum Levi--Civita connections on $(\mathbb Z_2)^4$}\\[0.4em]
\large A computer-assisted classification of a marked path--triangle sector}
\author{Marius Torsheim\\[0.3em]
\small Independent researcher\\
\small \AA{}saveien 14B, 0362 Oslo, Norway\\
\small\href{mailto:marius.torsheim123@gmail.com}{\texttt{marius.torsheim123@gmail.com}}}
\date{30 October 2026; manuscript revision 0.13}
\begin{document}\maketitle
\begin{abstract}
We define a marked, spatially homogeneous, trace-zero path--triangle
sector of bimodule quantum Levi--Civita connections for a specified
differential calculus on the four-dimensional binary cube. Reversible
algebraic reductions identify its real coefficient scheme with an explicit
localized family of four equations in four variables. We prove that its complexification
has generic geometric degree $512$. At the rational metric
$q^0=(4/3,12/11,3/4)$, with the four scales equal to one, we prove that
the entire declared fibre has length $512$ and consists of $512$ distinct
reduced points, exactly $204$ of which are real.
The proof separates the generic-degree calculation, the upper bound at
every finite special fibre, the construction of $512$ distinct lifts,
and the exact real-root count. Large integer polynomials and finite-field
identities are supplied as explicit computational certificates with
verification algorithms. The geometric arguments, exact identities and
standard algebraic-geometric ingredients are given explicitly. The count concerns the stated marked sector,
without a quotient by symmetry or a claim about the full moduli space.
\end{abstract}
\noindent\textbf{Keywords:} quantum Levi--Civita connections; finite geometry;
scheme-level reconstruction; computer-assisted enumeration; real algebraic geometry.

\section{Introduction}
Quantum Levi--Civita connections on a finite differential calculus are
specified by torsion, tensor metric compatibility and a star-preserving
bimodule Leibniz rule. Solving these equations is a nonlinear algebraic
problem even when the underlying algebra consists of functions on a
small finite set. An exact enumeration must retain the calculus, the
unknown tensor Leibniz map, the real structure and every open condition
through elimination and reconstruction.

We study the four coordinate-flip Grassmann calculus on
$(\mathbb Z/2)^4$, with a diagonal Lorentz-signature metric whose
coefficients are spatially homogeneous and may change between two
binary time slices. The family counted here is defined directly in the fixed marked connection data
by a spatial path on one slice, a spatial triangle on the other,
trace-zero conditions and an explicit regular open. These are coupling
graphs of the connection coefficients within one sixteen-vertex
calculus. The path is the minimally connected spatial support; the
triangle adds one independent cycle. The two slices are coupled by the
original temporal metric and star equations. The resulting completion
mechanisms and a balanced tensor-block splitting are explained in
Section~\ref{geometry:motivation}.

The contribution is a reversible description of this original localized
QLC scheme and a complete computer-assisted enumeration. Its generic
complex degree is $512$. At the rational metric
$q^0=(4/3,12/11,3/4)$, with all scales one, its entire declared fibre is
reduced with $512$ complex points, precisely $204$ of which are real.
The inverse reconstruction produces actual star-preserving bimodule
connections at all those real points. The theorem counts marked
coefficient points; its support and localization are retained throughout.
The explicit fibre supplies a test case for reconstruction and counting
methods that must retain both the original geometric equations and the
real structure while allowing the tensor Leibniz map to vary.

Three logically separate tasks are essential. The generic-degree
argument supplies the upper bound at every finite special fibre. Finiteness
and $512$ distinct smooth characteristic-zero lifts exhaust that bound.
Only then does an exact primitive-coordinate polynomial and a complete
real-root certificate identify the $204$ real connections. Finding $512$
modular points alone would leave the completeness question unanswered.
All written arguments and actual polynomial data are supplied. A single
supplementary command replays the arithmetic from fresh inputs; no earlier
project archive is required.

\subsection{Relation to prior work}
The finite-set bimodule framework, star-compatible connections and the
inner description of connections are established machinery
\cite{beggs-majid-star,majid-graphs}. In particular,
\cite[Definition~8.3, Proposition~8.11 and Lemma~8.17,
pp.~568, 576--579 and 589--590]{orig-book} gives the inner normal form,
torsion, full tensor metric compatibility and group component equations
used here. Its group-degree argument already forces the inner correction
map to vanish for the four coordinate flips. Star preservation is an
additional requirement, as in \cite[Lemma~8.7, p.~572]{orig-book}.
General component equations for nonconstant Cayley data also precede
this work \cite[Section~6.1, Proposition~6.2 and Lemma~6.3]{orig-graphs24},
\cite[Section~2.2 and Lemma~2.1]{orig-polygon}. The present result solves a specified
restricted coefficient scheme within that framework.

Square and polygon classifications already exhibit QLC nonuniqueness;
digital classifications also provide finite QLC counts
\cite{orig-square,orig-polygon,orig-digital}. The earlier generic square
family is already given in \cite[Section~5]{orig-emergence}. They depend
on the specified two-form calculus as well as the first-order graph.
For example, the cyclic and bit-flip squares have different two-form
translation degrees and torsion equations. In the bit-flip square,
$\deg(e^1\wedge e^2)=s_1+s_2\ne0$ in $(\mathbb Z/2)^2$;
in the cyclic four-gon,
$\deg(e^+\wedge e^-)=1+(-1)=0$ in $\mathbb Z/4$.
The basic two-form therefore has different bimodule commutation with
functions in these two calculi. Higher-dimensional component
analysis in \cite[Section~5.2, pp.~242--245]{orig-vh25} assumes cyclic orders $m_i>4$ and distinct
positive and negative directions. Setting those orders to two changes
the translation-degree blocks, so it is not a formal specialization of
that classification to the present input.

Uniqueness comparisons must keep their compatibility conventions and
fixed data. Theorem~5.14 of \cite{orig-mr25} compares connections with
the same tensor Leibniz map $\sigma$, whereas $\sigma$ varies here.
The inner formula below already implies uniqueness when our $\sigma$
is fixed. The weak Hermitian framework of \cite[Theorems~3.11 and~4.12]{orig-bgg} provides a
comparison of metric and torsion predicates on star-preserving invertible
bimodule pairs; it does not enumerate the marked support, traces and
localized scheme considered here. We make no claim to a new connection
formalism or a general uniqueness counterexample across different notions
of metric compatibility.

Further existence and uniqueness results require attention to the calculus.
Our one-form bimodule has zero centre: if $\omega=\sum_i a_ie^i$ is
central, then $a_i(f-R_if)=0$ for every vertex function $f$, and vertex
idempotents force all $a_i=0$. The centred-calculus hypothesis of
\cite[Definition~3.15 and Theorem~4.24]{orig-fmr26} fails here because
$Z(\Omega^1)=0$.  The derivation-based calculus of
\cite[Definitions~2.1 and~3.12, Theorem~7.2]{orig-ah26} has
$\mathfrak g\subseteq\operatorname{Der}(A)$, whereas
$\operatorname{Der}(\mathbb C^{16})=0$: an ordinary derivation
annihilates every primitive idempotent, since $D(p)=2pD(p)$ and $1-2p$
is a unit.
Aschieri's theorem \cite[Theorem~6.7]{orig-aschieri26} uses the dual of
the associated braided-derivation module.  On our natural Hopf input
$A=C(G)$, $G=(\mathbb Z/2)^4$, his cotriangular construction
\cite[Example~3.9, equations~(3.25) and~(3.28)]{orig-aschieri26}
gives zero one-forms: $U=A^*=\mathbb C[G]$ is commutative, so its adjoint
action is trivial and (3.28) reduces to the primitive-element equation,
whose only solution in $U$ is zero.  Thus this construction does not
give our rank-four Grassmann calculus.  This last comparison concerns
the natural Hopf input and makes no assertion about every possible
alternative triangular $H$-realization.

By contrast, \cite[Theorem~5.7]{orig-aw25} does apply to our bicovariant
calculus and the corresponding inverse metric
$\mathfrak g=\sum_i k_i\xi_i\otimes\xi_i$, $k_i=h_i^{-1}$.
Here $h_i$ are the metric coefficients of Section~\ref{geometry:model},
and $\xi_i$ is the right-linear dual frame with
$\langle\xi_i,e^j\rangle=\delta_{ij}$.
On symmetric-output coefficients $L^i{}_{jb}=L^i{}_{bj}$, their map
$\Phi_{\mathfrak g}$ sends $L$ to
$M_{ijb}=k_jL^i{}_{jb}+k_iL^j{}_{ib}$ and has inverse
$L^i{}_{jb}=(k_bM_{ijb}+k_jM_{ibj}-k_iM_{jbi})/(2k_jk_b)$.
Every denominator is a unit, so this proves invertibility here.
Using the notation of Section~\ref{geometry:model}, pass to a right
connection by $\nabla_R=D\circ\nabla\circ *$ and write
$\nabla_Re^i=\sum_{a,b}K^i_{ab}e^a\otimes e^b$.
Their compatibility condition
\cite[Definition~5.3(ii)]{orig-aw25} becomes
$k_jK^i_{jb}+k_iK^j_{ib}=\delta_{ij}(R_bk_i-k_i)$.
The case $i=j=b$, edge symmetry and
$K^i_{ii}=-\overline{U_i}$ force $U_i=0$, contradicting
\eqref{geometry:support}.  Moreover, pulling their unique connection
back by the same star passage and using \eqref{geometry:sigma} gives
the residual $E_{0ii}=-h_i(Q_i-1)^2/4$ in
\eqref{geometry:metriccomponents}, where
$Q_i(t)=h_i(1-t)/h_i(t)$ and $i>0$.  Thus the same-metric compatibility
conditions differ even without the marked trace restriction.
Both Aschieri and Aschieri--Weber allow arbitrary connections in these
theorems; these comparisons rely on their actual calculus and metric
predicates, rather than merely on fixing a tensor braiding.
Appendix~\ref{aw:appendix} derives the rational tensor-product
specialization, including the coefficient-shift cancellation, and the
same-metric residual explicitly.

Our contribution is the scheme-level reconstruction, the corrected
generic degree, the complete reduced rational-metric fibre and its exact
real count, taken together. The inspected versions are identified in the
bibliography. This statement does not assert a full four-cube moduli
classification, a general elimination algorithm, or universal priority.
Example~8.20 of \cite[pp.~594--598]{orig-book} and its relevant
dependencies have now been inspected in full. It treats the rank-two
bit-flip square: the rectangular metric family has four nonzero complex
connection parameters, and the displayed generic family has one; star
preservation imposes unit modulus. The intermediate case with only one
constant metric function is explicitly left without a separate
classification. Section~2.3, pp.~6--7 of the inspected arXiv v1 of
\cite{orig-square} has the same scope. These parameter families must not
be identified with the complexification of our restricted real scheme.
The general book machinery applies to our calculus and generates its
equations, but the inspected results do not supply the coupled marked
sector, its reversible reduction or the complete $512/204$ enumeration.
Section~\ref{geometry:motivation} also gives a concrete obstruction to
obtaining our connections by restriction to coordinate squares.
The detailed dependency and construction comparison is supplied in
\path{literature/literature_comparison.pdf}. The final twenty-three-page
journal version of \cite{orig-square} remains uninspected; no claim about
its complete classification scope is inferred from its metadata or from
the book's attribution of the example to that paper.

% BEGIN SHARED COUNTING EXPOSITION
\section{The statement and the proof obligations}
The geometry, all equations and the exact open defining the family are
given in Section~\ref{geometry:model}. The four scale parameters are
independent invertible base parameters; they are fixed in each counting
problem. The metric ratios vary over $q_i\ne0,1$.

\begin{theorem}[Degree and the exact rational-metric fibre]
\label{main:count}
Let $\mathcal Q_{\rm pt}$ be the real localized coefficient scheme defined
in Section~\ref{geometry:model}, and fix all four scales to one.
Let $X_{\mathbb C}$ be its complexification over the metric-ratio base.
Then:
\begin{enumerate}[label=\textup{(\roman*)}]
\item Its generic fibre has $512$ distinct geometric points and length
$512$. There is a nonempty base open on which the family is finite
\'etale of rank $512$.
\item At $q^0=(4/3,12/11,3/4)$, its entire declared fibre is finite,
reduced, and has $512$ distinct complex points and length $512$.
\item Exactly $204$ points in this fibre are real. Through the explicit
coefficient reconstruction, they are the star-preserving connections
in the specified real marked sector at this metric.
\end{enumerate}
\end{theorem}

The word ``complexification'' is essential: allowing the variables of
a real coefficient scheme to take complex values is not the same as
imposing an antilinear star equation on arbitrary complex-valued
connection coefficients. The $204$ real points have the latter intended
geometric interpretation through the real scheme isomorphism proved
below. Likewise positive metric parameters in the declared sign
convention do not mean a positive-definite coefficient matrix.

The proof is arranged in the logically necessary order
\[
 \begin{gathered}
 \text{generic degree }512
 \ \Longrightarrow\ \text{finite-fibre length cap }512\\
 \Longrightarrow\ \text{exactly }512\text{ points at }q^0
 \ \Longrightarrow\ 204\text{ real points}.
 \end{gathered}
\]
The first arrow is supplied by a separate length lemma. The second
also requires a proof that the special fibre is finite and a construction
of $512$ distinct characteristic-zero lifts. The last requires identifying
the true coordinate polynomial and proving a complete real-root count.
None of these steps is inferred merely from the number of modular
solutions found by a search.

This is a computer-assisted mathematical proof. Finite polynomial
certificates have the same role as long exact coefficient calculations
in an appendix: the data and specified arithmetic checks are part of
the proof, rather than numerical sampling. The arithmetic dependencies and exact replay scope are recorded in
Section~\ref{verification:section} and the computational supplement.

\section{The original connection problem and its exact algebraic model}
\label{geometry:model}

The sector counted below is specified here before any elimination or
specialization. Put
\[
 Q=q_1q_2q_3,\qquad
 R_{\mathbb Z}=\mathbb Z[q_1,q_2,q_3,
                 (2Q\textstyle\prod_i(q_i-1))^{-1}].
\]
For the connection problem use its real base change and four independent
metric scales:
\[
 K=(R_{\mathbb Z}\otimes_{\mathbb Z}\mathbb R)
       [\nu^{\pm1},\lambda_1^{\pm1},\lambda_2^{\pm1},\lambda_3^{\pm1}].
\]
The scales will subsequently all be fixed to one. They are independent
base parameters, not solution variables. Write
\[
 A_i(0)=\lambda_i,\quad A_i(1)=q_i\lambda_i,\qquad
 (h_0,h_1,h_2,h_3)(t)
       =(\nu,-A_1(t),-A_2(t),-A_3(t)).
\]
The positive metric parameter slice has
\(\nu,\lambda_i,q_i>0\). No positivity assumption is used in any of the
coordinate-ring identities proved in this section.

\subsection{Calculus, bimodule connection, metric and star equations}

Let \(B\) be an arbitrary commutative \(K\)-algebra; in particular it may
have nilpotents. Set \(B_{\mathbb C}=B[\mathrm i]/(\mathrm i^2+1)\), with
conjugation fixing \(B\), and let
\(\mathcal A=(B_{\mathbb C})^{(\mathbb Z/2)^4}\). Translation \(R_i\)
flips bit \(i\). The differential calculus and star are
\[
 e^if=R_i(f)e^i,\quad
 df=\sum_{i=0}^3(R_i f-f)e^i,\quad de^i=0,\qquad
 e^i\wedge e^j=-e^j\wedge e^i,\quad (e^i)^*=-e^i.
\]
All functions on the cube remain in \(\mathcal A\). Only the connection
and metric coefficients are spatially homogeneous: \(R_1,R_2,R_3\)
fix them, while \(R_0\) exchanges the two slices \(t=0,1\).
The diagonal tensor \(g=\sum_i h_i e^i\otimes e^i\) is central,
symmetric and real, and is nondegenerate because every \(h_i\) is a
unit. Its inverse pairing is
\((e^i,e^j)=\delta_{ij}R_i(h_i)^{-1}\).

A bimodule connection consists of \(\nabla\) and a bimodule map \(\sigma\)
with
\[
 \nabla(f\omega)=df\otimes\omega+f\nabla\omega,\qquad
 \nabla(\omega f)=(\nabla\omega)f+\sigma(\omega\otimes df).
\]
Here torsion freedom means \(\wedge\nabla=d\); metric compatibility
means \(\nabla g=0\), with the induced tensor connection; and star
preservation means \(\nabla(\omega^*)=\sigma D\nabla\omega\), where
\(D\) is conjugate tensor flip. No Yang--Baxter equation is imposed.

Let \(s_i\) be the four coordinate translations and
\(\theta=\sum_i e^i\), so \(df=\theta f-f\theta\).
The established inner normal form \cite[Proposition~8.11]{orig-book},
obtained from the two Leibniz rules, is
\[
 \nabla\omega=\theta\otimes\omega
                   -\sigma(\omega\otimes\theta)+\gamma(\omega),
\]
where \(\gamma\) is a bimodule map. A bimodule map preserves translation
degrees. The input degree of \(e^i\) has weight one, whereas every output
degree \(s_j+s_k\) has weight zero or two, so \(\gamma=0\).
This is valid over \(B\): vertex idempotents distinguish the translations
by coefficients \(0,1,-1\), without any reducedness hypothesis.

Use the sign convention
\(\nabla e^i=\sum_{j,k}C^i{}_{jk}e^j\otimes e^k\).
Translation degrees and torsion freedom give exactly
\begin{equation}\label{geometry:ansatz}
 C^i{}_{ii}=U_i,\qquad
 C^i{}_{ij}=C^i{}_{ji}=v_{ij},\qquad
 C^i{}_{jj}=w_{ij}\quad(j\ne i),
\end{equation}
with all three-distinct-index coefficients zero. There are \(56\)
complex coefficient values on both slices: \(8\) values of \(U\),
\(24\) of \(v\), and \(24\) of \(w\). The braiding is recovered uniquely:
\begin{equation}\label{geometry:sigma}
 \sigma^{il}{}_{jk}
  =\delta_{jl}\delta_{ki}
    -\mathbf1_{s_i+s_l=s_j+s_k}C^i{}_{jk}.
\end{equation}
Indeed, each allowed output has exactly one possible \(l\) in the inner
connection formula. Conversely this formula is a degree-preserving
bimodule map and supplies both Leibniz rules.

Expanding the tensor connection gives all metric components as
\begin{equation}\label{geometry:metriccomponents}
 E_{pqr}=\delta_{qr}(R_ph_r-h_r)+h_rC^r{}_{pq}
       +\sum_{i,j}h_iR_i(C^i{}_{jr})\sigma^{ij}{}_{pq}.
\end{equation}
The coefficient shifts follow from moving functions past \(e^i\).
Since the two minus signs in \((e^i)^*=-e^i\) cancel,
\[
 D(f e^j\otimes e^k)=R_jR_k(\bar f)e^k\otimes e^j.
\]
Thus the basis star components are
\begin{equation}\label{geometry:starcomponents}
 T^i{}_{pq}=C^i{}_{pq}
       +\sum_{j,k}R_jR_k(\overline{C^i{}_{jk}})\sigma^{kj}{}_{pq}.
\end{equation}
There are \(128\) metric and \(128\) star components on both slices.

Basis star equations are sufficient on all forms, as follows. Their
residual is
\[
 \nabla e^i+\sigma D\nabla e^i
       =\theta\otimes e^i-\sigma D\sigma(e^i\otimes\theta).
\]
For fixed \(i\), the inputs \(e^i\otimes e^l\) have distinct
translation degrees \(s_i+s_l\), preserved by \(D\) and \(\sigma\).
Consequently basis vanishing gives
\(\sigma D\sigma(e^i\otimes e^l)=e^l\otimes e^i\) separately for every
\(l\). Semilinearity then yields \((\sigma D)^2=\mathrm{id}\) on all
tensors, in particular \(\sigma^{-1}=D\sigma D\).
For
\[
 \mathcal R(\omega)=\nabla\omega-\sigma D\nabla(\omega^*)
\]
the two Leibniz rules, \((\sigma D)^2=\mathrm{id}\), and
\(D(\omega^*\otimes df^*)=df\otimes\omega\) show
\(\mathcal R(f\omega)=f\mathcal R(\omega)\).
Its vanishing on the basis therefore proves its vanishing on every form.
Applying this to \(\omega^*\) proves the stated star equation.

\subsection{The declared real scheme and its twenty-three units}

For each of the \(56\) coefficient values introduce independent real
coordinates \(c_R,c_I\), writing \(c=c_R+\mathrm i c_I\). Let \(I\)
be generated by the real and imaginary parts of all components
\eqref{geometry:metriccomponents}--\eqref{geometry:starcomponents}, and
of the following support and trace equations:
\begin{equation}\label{geometry:support}
 v_{13}(0)=v_{31}(0)=0,\qquad
 \sum_{i=0}^3(U_i(t)-1)=0\quad(t=0,1).
\end{equation}
There are \(112\) real coordinate variables before taking this quotient.
Define, before imposing the equations,
\[
 \begin{gathered}
 x_i(t)=v_{0i}(t),\quad y_i(t)=v_{i0}(t),\quad
 d_i(t)=x_i(t)+y_i(t)-1,\\
 X_{t,i}=x_i(1-t),\quad Y_{t,i}=y_i(t),\quad
 S_t=X_t^{\mathsf T}Y_t.
 \end{gathered}
\]
For two vectors write
\[
 E(X,Y)=X_1^2Y_2^2-X_1^2Y_3^2-X_2^2Y_1^2+X_2^2Y_3^2
                   +X_3^2Y_1^2-X_3^2Y_2^2,
\]
and put
\[
 A=X_{0,1}Y_{0,2}+X_{0,2}Y_{0,1},\quad
 B_0=X_{0,2}Y_{0,3}+X_{0,3}Y_{0,2},\quad E_1=E(X_1,Y_1).
\]
The notation \(B_0\) here is an edge pairing, not the test algebra \(B\).
Let \(\mathcal V\) contain both orientations of edges \(12,23\) on
slice zero and \(12,13,23\) on slice one; it has ten elements. Define
\begin{equation}\label{geometry:intrinsicopen}
 \Delta_{\rm reg}=
 \left(\prod_{i=1}^3d_i(0)(q_i-d_i(0))(q_i-d_i(0)^2)\right)
           S_0 A B_0 E_1\prod_{v\in\mathcal V}v,
 \qquad
 \mathcal Q=\Spec K[c_R,c_I,N(\Delta_{\rm reg})^{-1}]/I,
\end{equation}
where \(N(z)=z\bar z\). We also write
\(\mathcal Q_{\rm pt}=\mathcal Q\) to indicate this fixed marked
path--triangle sector. This is the precise regular open under study.
Its factors consist of nine temporal, four completion, and ten active-edge
factors, hence twenty-three in total.

For any test algebra, \(N(z)\) is a unit if and only if \(z\) is a unit
in \(B_{\mathbb C}\): if \(N(z)\) is invertible then
\(z^{-1}=\bar z/N(z)\); conversely conjugation preserves units.
The inverse of the real element $N(z)$ in $B_{\mathbb C}$ is fixed
by conjugation, so it lies in $B$; thus this equivalence also holds for
the norm as an element of $B$.
A product is invertible in a commutative ring if and only if each factor
is invertible. Thus \eqref{geometry:intrinsicopen} inverts exactly the
stated factors even before reality is proved. No pointwise active/inactive
dichotomy has been used.

\subsection{Eliminating the square coefficients and proving reality}

Antisymmetrizing \eqref{geometry:metriccomponents} in \(p,q\) cancels
the quadratic terms, because the lower connection indices are symmetric:
\begin{align}
 E_{pqr}-E_{qpr}
  ={}&\delta_{qr}(R_ph_r-h_r)-\delta_{pr}(R_qh_r-h_r)\notag\\
     &+h_qR_q(C^q{}_{pr})-h_pR_p(C^p{}_{qr}).
 \label{geometry:antisym}
\end{align}
For the six unordered pairs, \(r=p,q\), and the two slices these are
twenty-four linear equations with distinct \(w\)-pivots
\(\pm\nu\) or \(\pm A_i(t)\). Their unique solution is
\begin{align}
 w_{i0}(t)&=-\nu x_i(1-t)/A_i(t),\notag\\
 w_{0i}(t)&=\{A_i(1-t)(1-y_i(1-t))-A_i(t)\}/\nu,\notag\\
 w_{ij}(t)&=A_j(t)v_{ji}(t)/A_i(t)\quad(i,j\text{ spatial},\ i\ne j).
 \label{geometry:wreconstruction}
\end{align}
All pivots are units, so this is reversible elimination of variables in
coordinate rings, without a radical operation.

For the next formulas suppress \(t\), use a tilde for the opposite
slice, and put
\[
 q_i(t)=A_i(1-t)/A_i(t),\quad m_i=U_i-1,\quad
 M_{ii}=m_i,\ M_{ij}=v_{ij}\ (i\ne j\text{ spatial}),\quad
 z_i=1+q_i(t)(\widetilde y_i-1),
\]
\[
 \widehat L=\begin{pmatrix}m_0&X^{\mathsf T}\\z&M\end{pmatrix},\qquad
 L^\sharp=\begin{pmatrix}\widetilde m_0&X^{\mathsf T}\\y&M\end{pmatrix}.
\]
Substitution of \eqref{geometry:wreconstruction} into the original
components gives the following exhaustive table. Counts include both
slices and duplicated components.
\begin{center}\small
\begin{tabular}{p{.31\linewidth}rp{.55\linewidth}}
\hline
Indices&Count&Exact component after substitution\\\hline
\(p=q\)&32&
\(E_{00r}=-\nu(\widehat LL^\sharp-\Id)_{0r}\),
\(E_{iir}=A_i(t)(\widehat LL^\sharp-\Id)_{ir}\).\\
\(p\ne q\), time--space, \(r=p\) or \(q\)&24&
\(-\nu(d_i\widetilde x_i-x_i)\) for \(r=0\);
\(A_i(t)(d_i^2-d_ix_i+q_i(t)(x_i-1))\) for \(r=i\).\\
\(p\ne q\), spatial, \(r=p\) or \(q\)&24&
\(A_r(t)v_{rj}(v_{pq}+v_{qp}-2)\), where \(j\) is the other index.\\
Three distinct indices&48&Zero.\\\hline
\end{tabular}
\end{center}
For example the time--space entries follow by inserting the two displayed
\(w\)-formulas into \eqref{geometry:metriccomponents}; the remaining
off-diagonal terms have the single possible spatial mixed coefficient.
For diagonal outputs the four sums are precisely the matrix product shown.
These identities cover every component and prove equality of the original
metric ideal and the displayed residual ideal after the reversible linear
elimination, rather than just equality of field-valued zero sets.

For the star table define \(W_{ij}=C^i{}_{jj}\),
\(\mathscr Q=\Id-W\), and \(d_{ij}=v_{ij}+v_{ji}-1\).
Direct use of \eqref{geometry:sigma} gives
\begin{center}\small
\begin{tabular}{p{.31\linewidth}rp{.55\linewidth}}
\hline
Output indices&Count&Exact star component\\\hline
\(p=q\)&32&\((\Id-\overline{\mathscr Q}\mathscr Q)_{ip}\).\\
\(p\ne q\), \(i\in\{p,q\}\)&48&
\(v_{ij}-d_{ij}R_iR_j(\bar v_{ij})\), with \(j\) the other index.\\
Three distinct indices&48&Zero.\\\hline
\end{tabular}
\end{center}
Summing mixed orientations yields
\(1-d_{ij}R_iR_j(\bar d_{ij})\). Thus all \(d_{ij}\) are units.
Also \(\overline{\mathscr Q}\mathscr Q=\Id\) makes
\(\mathscr Q\) invertible by the adjugate identity. The degree-zero
braiding block is \(\mathscr Q^{\mathsf T}\), and the mixed block for
\(i,j\) has matrix
\(\left(\begin{smallmatrix}-v_{ij}&1-v_{ji}\\
1-v_{ij}&-v_{ji}\end{smallmatrix}\right)\)
and determinant \(d_{ij}\). This also checks invertibility of the entire
braiding directly.

Now consider one time--space pair on slice zero and abbreviate
\[
 x=x_i(0),\quad x'=x_i(1),\quad y=y_i(0),\quad y'=y_i(1),
 \quad d=x+y-1,\quad d'=x'+y'-1,\quad q=q_i.
\]
The metric residuals include
\[
 F_d=dx'-x,\qquad F'_d=d'x-x',\qquad
 H_d=d^2-dx+q(x-1),\qquad (q-d)x=q-d^2.
\]
Our open makes \(d,q-d,q-d^2\) units. Hence \(x\) and \(x'=x/d\)
are units, and the polynomial identity
\[
 dF'_d+F_d=x(dd'-1)
\]
gives \(dd'=1\). The star residuals are
\(G=x-d\bar x'\), \(G'=x'-d'\bar x\). The identities
\[
 G+F_d=d(x'-\bar x'),\qquad
 G'+F'_d=d'(x-\bar x)
\]
prove \(x,x'\) real in the universal localized coordinate ring.
Consequently \(d=x/x'\), \(d'=d^{-1}\), \(y=d+1-x\), and
\(y'=d'+1-x'\) are real as well. This is a statement in the ring:
\(B_{\mathbb C}=B\oplus\mathrm iB\), so
\(c-\bar c=2\mathrm i c_I=0\) forces \(c_I=0\), even for nilpotent
imaginary coordinates.

The temporal entries have therefore been recovered uniquely as
\begin{equation}\label{geometry:temporal}
 x=\frac{q-d^2}{q-d},\quad
 y=\frac{d(q-1)}{q-d},\quad x'=x/d,\quad y'=y/d.
\end{equation}
All four are units. They satisfy the opposite-slice conic with
\((q,d)\) replaced by \((q^{-1},d^{-1})\).
The replacement \(z=y\), which cannot be made before this step, follows
from the polynomial certificate
\[
 d\{1+q(y'-1)-y\}=-H_d+q(dd'-1-F_d).
\]
Thus the metric core becomes \(LL^\sharp=\Id\), with
\[
 L=\begin{pmatrix}m_0&X^{\mathsf T}\\y&M\end{pmatrix},\qquad
 L^\sharp=\begin{pmatrix}\widetilde m_0&X^{\mathsf T}\\y&M\end{pmatrix}.
\]
A right inverse of a square matrix over a commutative ring is also a left
inverse: its determinant is a unit, and the adjugate gives the inverse.
The upper-right block of \(LL^\sharp-L^\sharp L\) is
\((m_0-\widetilde m_0)X^{\mathsf T}\). Every \(X_i=x'_i\) is a unit,
so \(m_0=\widetilde m_0=:m\). On each slice put \(P_t=L=L^\sharp\);
then \(P_t^2=\Id\).

On each active spatial edge one orientation is a declared unit. The metric
product in the table consequently gives \(v_{ij}+v_{ji}=2\). On the
inactive path edge both orientations are zero by
\eqref{geometry:support}. With
\(H_t=\operatorname{diag}(\nu,-A_1(t),-A_2(t),-A_3(t))\), the
reconstruction \eqref{geometry:wreconstruction} gives the exact identity
\[
 \mathscr Q=-H_t^{-1}P_t^{\mathsf T}H_t.
\]
The metric scales are real. Therefore the degree-zero star equations give
\(P_t\bar P_t=\Id\). Combined with \(P_t^2=\Id\), they imply
\(\bar P_t=P_t\), with the explicit polynomial certificate
\[
 \bar P-P=P(P\bar P-\Id)-(P^2-\Id)\bar P.
\]
All \(U,v\), and then all \(w\), are real in the coordinate ring.
The trace equations are exactly \(\tr P_t=0\).

\subsection{The temporal chart and the block problem}

Write \(u_i=y_i(0)\), and set
\[
 \alpha_i=q_i-(u_i-1)^2,\qquad \beta_i=q_i+u_i-1.
\]
The inverse changes of coordinates, together with every required unit
identity, are
\begin{align}
 u_i&=\frac{d_i(q_i-1)}{q_i-d_i},&
 \beta_i&=\frac{q_i(q_i-1)}{q_i-d_i},&
 \alpha_i&=\frac{q_i(q_i-1)(q_i-d_i^2)}{(q_i-d_i)^2},\notag\\
 d_i&=\frac{q_iu_i}{\beta_i},&
 q_i-d_i&=\frac{q_i(q_i-1)}{\beta_i},&
 q_i-d_i^2&=\frac{q_i(q_i-1)\alpha_i}{\beta_i^2}.
 \label{geometry:chartinverse}
\end{align}
They prove isomorphism of the two temporal localizations, including
nonreduced test algebras. The two slice vectors become exactly
\begin{equation}\label{geometry:XY}
 X_{0,i}=\frac{\alpha_i}{q_iu_i},\quad Y_{0,i}=u_i,
 \qquad
 X_{1,i}=\frac{\alpha_i}{\beta_i},\quad Y_{1,i}=\frac{\beta_i}{q_i}.
\end{equation}
Every entry is a unit, and
\[
 z_i=X_{t,i}Y_{t,i}=\alpha_i/q_i,\qquad S=S_0=S_1=\sum_i z_i.
\]
Thus the intermediate block problem has
\[
 P_t=\begin{pmatrix}m&X_t^{\mathsf T}\\Y_t&M_t\end{pmatrix},\qquad
 M_0=\begin{pmatrix}a_0&p_0&0\\2-p_0&b_0&s_0\\0&2-s_0&c_0\end{pmatrix},
\quad
 M_1=\begin{pmatrix}a_1&p_1&r_1\\2-p_1&b_1&s_1\\2-r_1&2-s_1&c_1\end{pmatrix},
\]
subject to \(P_t^2=\Id\), \(\tr P_t=0\), and precisely the
twenty-three units already declared. The subscripts on \(a_t,b_t,c_t\)
are slice labels, unrelated to the metric scales.

\begin{lemma}\label{geometry:blockcriterion}
Over any commutative ring with \(2,S=X^{\mathsf T}Y\) invertible, suppose
\(MY=-mY\), \(X^{\mathsf T}M=-mX^{\mathsf T}\), and
\(\tr M=-m\). Then the block matrix
\(P=\left(\begin{smallmatrix}m&X^{\mathsf T}\\Y&M\end{smallmatrix}\right)\)
is an involution if and only if
\[
 F=m^2+S-1=0,\qquad G=\tr(M^2)-m^2-2=0.
\]
This is equality of equation ideals after the linear and trace equations.
\end{lemma}
\begin{proof}
Put \(\Pi=YX^{\mathsf T}/S\). The free module of rank three splits
as \(AY\oplus\ker X^{\mathsf T}\). The second summand is a projective
module of rank two; the linear equations make \(M\) act by \(-m\) on
the first summand and by a trace-zero endomorphism \(N\) on the second.
Locally a trace-zero matrix
\(\left(\begin{smallmatrix}a&b\\c&-a\end{smallmatrix}\right)\)
squares to \((a^2+bc)\Id\), half the trace of its square.
These endomorphism identities glue. It follows that
\begin{equation}\label{geometry:blockresidual}
 M^2+YX^{\mathsf T}-\Id=\frac G2(\Id-\Pi)+F\Pi.
\end{equation}
The other blocks of \(P^2-\Id\) are \(F\) and the assumed linear
residuals. Conversely, if \(R_M=M^2+YX^{\mathsf T}-\Id\), then
\(G=\tr R_M-F\). Both inclusions of ideals follow in the universal
quotient by the linear and trace equations.
\end{proof}

The localization at \(S\) is substantive. For example
\(m=1\), \(X=(1,1,1)^{\mathsf T}\), \(Y=(1,1,-2)^{\mathsf T}\),
and
\(M=\left(\begin{smallmatrix}0&-1&0\\-1&0&0\\0&0&-1\end{smallmatrix}\right)\)
have \(S=0\), the linear and trace equations, and \(F=G=0\), but
\(M^2+YX^{\mathsf T}-\Id=YX^{\mathsf T}\ne0\).
This tests the lemma, without asserting that the example is in our
selected support sector.

\subsection{Path completion and its three scalar equations}

For this subsection \(X,Y\) are the path vectors. Abbreviate
\[
 A=X_1Y_2+X_2Y_1,\quad B=X_2Y_3+X_3Y_2,\quad z_i=X_iY_i,
\]
\[
 C=A(z_2+z_3)+B(z_1+z_2),\qquad
 L=ABm+C,\qquad H_p=AB+4z_2.
\]
The six linear equations \(M_0Y=-mY\),
\(X^{\mathsf T}M_0=-mX^{\mathsf T}\) have the unique solution
\begin{align}
 p_0&=\frac{2X_2Y_1}{A},&s_0&=\frac{2X_3Y_2}{B},\notag\\
 a_0&=-m-\frac{2z_2}{A},&c_0&=-m-\frac{2z_2}{B},&
 b_0&=-m-\frac{2z_1}{A}-\frac{2z_3}{B}.
 \label{geometry:pathcompletion}
\end{align}
For uniqueness, take \(X_1\) times the first right residual minus
\(Y_1\) times the first left residual. It is
\(Ap_0-2X_2Y_1\). The analogous residual at vertex three gives
\(Bs_0-2X_3Y_2\). Since \(A,B,Y_i\) are units, these determine
the two edges and all diagonal entries. Direct substitution verifies the
remaining linear residuals.

After this completion put \(F_1=\tr M_0+m\) and
\(F_2=\tr(M_0^2)-m^2-2\). Expansion gives
\begin{align}
 F_1&=-2L/(AB),\notag\\
 AB(F_2-2F)+2CF_1+2SH_p&=0.
 \label{geometry:pathsyzygy}
\end{align}
For example \(p_0(2-p_0)=4z_1z_2/A^2\) and
\(s_0(2-s_0)=4z_2z_3/B^2\); these and the three diagonal squares
give the second identity. Explicit inverse generator transformations are
\[
 F_2=2F+\frac{4CL}{(AB)^2}-\frac{2SH_p}{AB},\qquad
 H_p=-\frac{AB(F_2-2F)+2CF_1}{2S}.
\]
Every denominator is a declared unit. Hence
\((F,F_1,F_2)=(F,L,H_p)\) as ideals, and the trace-zero path
involution is equivalent to \(F=L=H_p=0\).

\subsection{Triangle completion and an explicit invariant residual}

Now \(X,Y\) are the triangle vectors. The unknowns
\((a_1,b_1,c_1,p_1,r_1,s_1)\) are determined by the three right
linear equations, the first two left equations, and the trace. Their matrix
and right-hand side are
\[
 \mathcal T=\begin{pmatrix}
 Y_1&0&0&Y_2&Y_3&0\\
 0&Y_2&0&-Y_1&0&Y_3\\
 0&0&Y_3&0&-Y_1&-Y_2\\
 X_1&0&0&-X_2&-X_3&0\\
 0&X_2&0&X_1&0&-X_3\\
 1&1&1&0&0&0
 \end{pmatrix},\qquad
 \mathbf b=\begin{pmatrix}
 -mY_1\\-2Y_1-mY_2\\-2Y_1-2Y_2-mY_3\\
 -mX_1-2X_2-2X_3\\-mX_2-2X_3\\-m
 \end{pmatrix}.
\]
Expansion of its determinant gives \(\det\mathcal T=Y_3E(X,Y)\),
which is a unit on our open. Thus the solution is unique over every test
ring. The omitted third left equation follows from
\[
 X^{\mathsf T}(MY+mY)=(X^{\mathsf T}M+mX^{\mathsf T})Y
\]
and the unit \(Y_3\).
Define polynomial numerators, in the displayed coordinate order, by
\[
 (\Theta_a,\Theta_b,\Theta_c,\Theta_p,\Theta_r,\Theta_s)^{\mathsf T}
       =\operatorname{adj}(\mathcal T)\mathbf b/Y_3.
\]
Each component of the numerator is divisible by \(Y_3\) in the integer
polynomial ring; this can equivalently be checked by the determinant
expansion just specified. Thus this is an explicit polynomial definition,
not a choice of a rational solution. The completed matrix is \(M_1=\mathcal N/E\),
where
\[
 \mathcal N=\begin{pmatrix}
 \Theta_a&\Theta_p&\Theta_r\\
 2E-\Theta_p&\Theta_b&\Theta_s\\
 2E-\Theta_r&2E-\Theta_s&\Theta_c
 \end{pmatrix}.
\]
Let \(G_{\rm num}=\tr(\mathcal N^2)-(m^2+2)E^2\) and
\(Q_\triangle=[m^2]G_{\rm num}\). One has the polynomial identity
\begin{equation}\label{geometry:trianglesyzygy}
 G_{\rm num}=Q_\triangle F-2S H_{\rm tri}(X,Y,m).
\end{equation}
Here \(Q_\triangle\) is independent of \(m\), and the following
orbit formulas give an explicit residual \(H_{\rm tri}=H_0+mH_1\).
Write \(z_i=X_iY_i\), \(a_{ij}=X_iY_j+X_jY_i\). For a monomial
\(f\), \(\Orb(f)\) denotes the sum of its distinct simultaneous
\(S_3\)-images, permuting vertex and unordered-edge labels. Set
\begin{align}
 H_1={}&16\Orb(z_1z_2a_{12})+8\Orb(z_1a_{12}a_{23})
          -4\Orb(a_{12}^3)+4\Orb(a_{12}^2a_{13}),\notag\\
 H_0={}&16\Orb(z_1^2z_2)+16\Orb(z_1^2a_{23})+32z_1z_2z_3\notag\\
 &-4\Orb(z_1z_2a_{12}^2)+16\Orb(z_1z_2a_{13})
       -4\Orb(z_1a_{12}a_{13}a_{23})\notag\\
 &+8\Orb(z_1a_{12}a_{23})+\Orb(a_{12}^4)+4\Orb(a_{12}^3)\notag\\
 &-2\Orb(a_{12}^2a_{13}^2)-4\Orb(a_{12}^2a_{13})
       +16a_{12}a_{13}a_{23}.
 \label{geometry:invariants}
\end{align}
The identity \eqref{geometry:trianglesyzygy} is an identity after the
bilinear substitution, with no algebraically independent extra bilinear
variables. In particular
\[
 4z_1z_2z_3+a_{12}a_{13}a_{23}
        -z_1a_{23}^2-z_2a_{13}^2-z_3a_{12}^2=0
\]
holds identically. Equations \eqref{geometry:trianglesyzygy} and
\eqref{geometry:invariants} are directly verified by replacing each
\(\Theta\) with the adjugate definition and expanding; they are finite
integer polynomial identities. The accompanying executable certificate
performs exactly this expansion, including all six numerators, all seven
linear/trace constraints, and the displayed orbit polynomial.

Since \(2,S,E\) are units,
\[
 (F,\tr(M_1^2)-m^2-2)=(F,H_{\rm tri})
\]
after completion. Indeed
\(E^2(\tr(M_1^2)-m^2-2)=Q_\triangle F-2SH_{\rm tri}\)
exhibits both generator transformations. Lemma
\ref{geometry:blockcriterion} now gives equivalence with the full
trace-zero triangle involution.

\subsection{The four-variable family, its cleared equations and exact open}

Use \eqref{geometry:XY}, and use \(A,B\) for the path pairings from
now on. The resulting family has four coordinates \((u_1,u_2,u_3,m)\)
and exactly four rational equations:
\begin{equation}\label{geometry:four}
 \boxed{\begin{aligned}
 H_p&=AB+4z_2=0,\\
 L&=ABm+A(z_2+z_3)+B(z_1+z_2)=0,\\
 F&=m^2+\sum_i z_i-1=0,\\
 H_{\rm tri}&=H_0(z,a_1)+mH_1(z,a_1)=0.
 \end{aligned}}
\end{equation}
In the last equation \((a_1)_{ij}=X_{1,i}Y_{1,j}+X_{1,j}Y_{1,i}\).
Denote this localized scheme over \(R_{\mathbb Z}\) by \(\mathcal X\).
The open is intrinsically defined by the temporal, completion and active
edge units above. The next formulas give its exact polynomial inventory.

Put
\[
 \Ab=q_2\alpha_1u_2^2+q_1\alpha_2u_1^2,\qquad
 \Bb=q_3\alpha_2u_3^2+q_2\alpha_3u_2^2,\qquad
 \overline S=Q\sum_i\alpha_i/q_i.
\]
Then
\(A=\Ab/(q_1q_2u_1u_2)\) and
\(B=\Bb/(q_2q_3u_2u_3)\). The full inventory is
\begin{center}\small
\begin{tabular}{p{.29\linewidth}p{.55\linewidth}r}
\hline
Intrinsic units&Cleared polynomial factors&Count\\\hline
Temporal chart&\(u_i,\alpha_i,\beta_i\ (i=1,2,3)\)&9\\
\(S\)&\(\overline S\)&1\\
\(A,B\)&\(\Ab,\Bb\)&2\\
Four path entries&\(q_1\alpha_2u_1^2,\ q_2\alpha_1u_2^2,
q_2\alpha_3u_2^2,\ q_3\alpha_2u_3^2\)&4\\
\(E_1\), six triangle entries&The seven polynomials defined immediately below&7\\\hline
\end{tabular}
\end{center}
For example
\[
 p_0=2q_1\alpha_2u_1^2/\Ab,\qquad
 2-p_0=2q_2\alpha_1u_2^2/\Ab,
\]
and the other two path entries have the analogous formulas with \(\Bb\).
Some listed factors are redundant on this open. They are retained, so no
change of sector is hidden in simplification.

For a triangle polynomial \(J=\sum c_{x,y,k}X^xY^ym^k\), define
\[
 b_i(J)=\max_{c_{x,y,k}\ne0}(x_i-y_i),\quad
 c_i(J)=\max_{c_{x,y,k}\ne0}y_i,\qquad
 \mathcal C_1(J)=\left(\prod_i\beta_i^{b_i(J)}q_i^{c_i(J)}\right)
                      J(X_1,Y_1,m).
\]
Termwise use of \eqref{geometry:XY} shows that this is a polynomial:
the exponents are \(b_i+y_i-x_i\ge0\) and \(c_i-y_i\ge0\).
All clearing monomials are units on the temporal chart. The remaining
seven factors are precisely
\[
 \begin{gathered}
 \mathcal C_1(E),\quad
 \tfrac12\mathcal C_1(\Theta_p),\quad
 \tfrac12\mathcal C_1(2E-\Theta_p),\quad
 -\tfrac12\mathcal C_1(\Theta_r),\quad
 \tfrac12\mathcal C_1(2E-\Theta_r),\quad
 \tfrac12\mathcal C_1(\Theta_s),\quad
 -\tfrac12\mathcal C_1(2E-\Theta_s).
 \end{gathered}
\]
The constant signs and factors \(2\) are units. Since the six completed
triangle entries are \(\Theta/E\) or \((2E-\Theta)/E\), the polynomial
and intrinsic opens agree over every test algebra.

For a polynomial rather than rational presentation, adjoin five graph
coordinates \(\xi,\eta,c_{12},c_{13},c_{23}\). Their defining equations
are
\begin{align}
 q_1q_2u_1u_2\xi-\Ab&=0,&
 q_2q_3u_2u_3\eta-\Bb&=0,\notag\\
 q_iq_j\beta_i\beta_j
 (c_{ij}-X_{1,i}Y_{1,j}-X_{1,j}Y_{1,i})&=0
 &&(ij=12,13,23).
 \label{geometry:graph}
\end{align}
Their coefficients on the new variables are units, so these variables
adjoin unique regular sections. The other four polynomial generators are
\begin{equation}\label{geometry:cleared}
 \begin{split}
 &q_2(\xi\eta+4\alpha_2/q_2),\\
 &Q\{\xi\eta m+\xi(z_2+z_3)+\eta(z_1+z_2)\},\\
 &Q(m^2+z_1+z_2+z_3-1),\\
 &-Q^4\{H_0(z,c)+mH_1(z,c)\}.
 \end{split}
\end{equation}
They are integer polynomials; the base units clear all denominators.
Together \eqref{geometry:graph} and \eqref{geometry:cleared} give the
nine-generator presentation of the same localized four-variable family.
The complete coefficients of these generators and of all twenty-three
factors are supplied, and the symbolic certificate checks the displayed
definitions against every coefficient.

\subsection{Binding the ten-equation source and the centre equations}
\label{geometry:centre-source}

An alternative polynomial presentation keeps the six triangle entries
\((a,b,c,p,r,t)=(a_1,b_1,c_1,p_1,r_1,s_1)\). Put
\[
 M=\begin{pmatrix}a&p&r\\2-p&b&t\\2-r&2-t&c\end{pmatrix},\qquad
 \mathcal K=\Ab(q_3\alpha_2+q_2\alpha_3)u_3
             +\Bb(q_2\alpha_1+q_1\alpha_2)u_1.
\]
The ten literal source equations used in the exact computation are the
following polynomials, in the coordinate order
\((u_1,u_2,u_3,m,a,b,c,p,r,t)\):
\begin{align}
 &Q(m^2-1)+\sum_iq_{i+1}q_{i+2}\alpha_i,\notag\\
 &\Ab\Bb+4Q\alpha_2u_1u_2^2u_3,\qquad
 m\Ab\Bb+u_2\mathcal K,\notag\\
 &\sum_{j=1}^3(M_{ij}+m\delta_{ij})\beta_jq_{j+1}q_{j+2}
                    \quad(i=1,2,3),\notag\\
 &\sum_{j=1}^3(M_{ji}+m\delta_{ij})\alpha_j\beta_{j+1}\beta_{j+2}
                    \quad(i=1,2),\notag\\
 &a+b+c+m,\notag\\
 &a^2+b^2+c^2+2\{p(2-p)+r(2-r)+t(2-t)\}-m^2-2.
 \label{geometry:tensource}
\end{align}
The indices on \(q,\alpha,\beta\) are cyclic modulo three. To verify
scheme equivalence, set
\(D_p=q_1q_2^2q_3u_1u_2^2u_3\), a temporal unit.
The first three source rows are exactly \(QF,D_pH_p,D_pL\).
The next three are \(Q(MY_1+mY_1)_i\); the next two are
\((\beta_1\beta_2\beta_3)(X_1^{\mathsf T}M+mX_1^{\mathsf T})_i\).
The last two are the trace and trace-square residuals. All multipliers
are units. The triangle linear system is invertible on our declared
\(Y_{1,3}E_1\) open. The omitted left equation and the inverse
transformations in \eqref{geometry:trianglesyzygy} therefore prove
equivalence with \eqref{geometry:four} as localized schemes. A separate
exact integer-polynomial comparison checks every literal coefficient of
\eqref{geometry:tensource}; it is not a modular comparison.

One must also specify the representative used to eliminate \(m\). Since
\(z_2=\alpha_2/q_2\) is a unit, define
\[
 m_* = C/(4z_2).
\]
The exact identity
\begin{equation}\label{geometry:melimination}
 L=AB(m-m_*)+m_*H_p
\end{equation}
gives \((H_p,L)=(H_p,m-m_*)\) on the path-completion open.
The expression \(-C/(AB)\), obtained from \(L=0\) alone, agrees with
\(m_*\) modulo \(H_p\); they are not identical rational functions.
Indeed their difference is \(CH_p/(4z_2AB)\).

At \(q^0=(4/3,12/11,3/4)\), eliminating \(m=m_*\) gives the three
primitive integer centre numerators \(P_p,P_F,P_T\), of respective
\(u_1\)-degrees \(2,6,10\). Their canonical scalar identities are
\eqref{centre:source1}--\eqref{centre:source3}. Every multiplier is a
declared unit. The binding checker reconstructs these rational residuals
over \(\mathbb Q[u_1,u_2,u_3]\), checks zero remainder after division by
the complete supplied polynomial, and verifies all multiplier factors.
Thus these are equalities of localized ideals over \(\mathbb Q\).

\subsection{Inverse reconstruction and the scheme-theoretic conclusion}

\begin{theorem}\label{geometry:bridge}
The real scheme \(\mathcal Q\) of
\eqref{geometry:intrinsicopen} is isomorphic over \(\Spec K\) to
\(\mathcal X\times_{\Spec R_{\mathbb Z}}\Spec K\).
All imaginary connection coordinates vanish in its coordinate ring.
The isomorphism is natural over arbitrary commutative \(K\)-algebras,
including nonreduced algebras. It identifies actual real star-preserving
connections with the real points of \(\mathcal X\); its complexification
is the complexification of this real algebraic model.
\end{theorem}
\begin{proof}
The forward construction has already been carried out in the universal
localized coordinate ring. It reads
\(d_i=x_i(0)+y_i(0)-1\), \(u_i=y_i(0)\), \(m=U_0(0)-1\), and
the spatial \(M_t\) from the original coefficients. The component tables
and unit pivots prove reality, common \(m\), the two block involutions
and their traces. The temporal coordinate inverse, unique path and triangle
completions, and the two residual identities prove all four equations
\eqref{geometry:four}. The explicit factor inventory proves that this
map lands in the exact declared open.

Conversely, start with an \(\mathcal X\)-point over a \(K\)-algebra.
Set \(d_i=q_iu_i/\beta_i\), and reconstruct all time--space entries
by \eqref{geometry:temporal}; equivalently use the two vectors
\eqref{geometry:XY}. Construct \(M_0\) by
\eqref{geometry:pathcompletion} and \(M_1=\mathcal N/E\) by the
adjugate formula. The inverse ideal identities in
\eqref{geometry:pathsyzygy} and \eqref{geometry:trianglesyzygy}, with
Lemma \ref{geometry:blockcriterion}, give both trace-zero block
involutions. Set
\[
 U_0(0)=U_0(1)=1+m,\qquad
 U_i(t)=1+(M_t)_{ii},\qquad v_{ij}(t)=(M_t)_{ij}
       \quad(i\ne j\text{ spatial}),
\]
recover all \(w\) using \eqref{geometry:wreconstruction}, and set all
imaginary coordinates to zero. This reconstructs all \(56\) complex
coefficient values without any root choice.

The lower-index symmetry verifies torsion freedom. Formula
\eqref{geometry:sigma} and the inner connection formula give the
bimodule connection. The temporal identities verify all time--space metric
residuals; the active edge sums equal two and the inactive pair is zero;
the block involutions verify every metric-core residual. The exhaustive
metric table therefore proves \(E_{pqr}=0\) for every component.

For star preservation, on an active spatial edge \(d_{ij}=1\) and the
entries are real; on the inactive edge both entries vanish. On a temporal
edge \eqref{geometry:temporal} gives \(x=dx'\), \(y=dy'\) and their
opposite-slice counterparts. Finally
\(\mathscr Q=-H_t^{-1}P_t^{\mathsf T}H_t\) is real with
\(\mathscr Q^2=\Id\), so the zero-degree star block vanishes.
The exhaustive star table verifies every \(T^i{}_{pq}=0\); the argument
following \eqref{geometry:starcomponents} extends this to all forms and
proves invertibility of \(\sigma\). All original support, trace, and
regular-open conditions are satisfied.

The forward and backward formulas are inverse: the temporal coordinate
maps are inverse by \eqref{geometry:chartinverse}; the path and triangle
entries are unique solutions with unit pivots; and the eliminated square
coefficients have unique unit-leading linear equations. They commute with
arbitrary maps of test algebras. They thus define inverse maps of localized
coordinate rings, proving the asserted scheme isomorphism, not merely a
bijection of classical points.
\end{proof}

In particular, after all four scales are fixed to one and
\(q^0=(4/3,12/11,3/4)\), a real point counted later is an actual
star-preserving connection in this precise sector. The complex count is
the geometric count in the complexification of its real coordinate scheme.
It does not assert that arbitrary complex points, with an independently
imposed analytic conjugation operation, are star-preserving connections.
The omitted exceptional parameters, other supports, other traces,
completion-degenerate charts, and unmarked quotients are not part of the
defined counting problem.

\section{Preservation of the original QLC problem and the meaning of the sector}
\label{integrity:section}

Theorem~\ref{geometry:bridge} gives the scheme isomorphism. The following
tables record how its two maps transport the geometric requirements:
fixed data and sector restrictions, all inverted factors, and the original
QLC residuals. The final subsection explains the geometric meaning of the
chosen support and trace. The bimodule conventions are those of
\cite{beggs-majid-star,majid-graphs}.

\subsection{Fixed data, imposed restrictions and derived consequences}
\label{integrity:choices}

The base space has all sixteen vertices of $(\mathbb Z/2)^4$ and the
four specified translation directions throughout. Spatial homogeneity
restricts coefficients, not the algebra of functions or the differential
calculus. The ordered frame and the two labelled slices are retained.
The following distinctions prevent a consequence of the equations from
being silently substituted for an additional sector assumption.

\begin{center}\small
\begin{tabular}{p{.23\linewidth}p{.32\linewidth}p{.36\linewidth}}
\hline
Item&Status before reduction&What is established afterwards\\\hline
Calculus and metric&Fixed as in Section~\ref{geometry:model}; the four
scales and $q_i$ are base parameters.&All bimodule, metric and star
calculations use the same shifts and sign convention.\\
Spatial homogeneity&Imposed on coefficients on each slice.&The resulting
components still define connections on the full sixteen-point algebra.\\
Spatial support&Both $13$ orientations vanish on slice zero; the ten
specified orientations are units.&Exactly the path and triangle
coupling patterns remain; active pairs have sum two.\\
Trace&$\sum_i(U_i(t)-1)=0$ is imposed on each slice.&It becomes
$\tr P_t=0$; it is not deduced from torsion or metric compatibility.\\
Reality and star&The original star equations are imposed in real and
imaginary coordinates.&Every imaginary coefficient vanishes in the
localized coordinate ring; this is proved, not assumed.\\
Regular open&The twenty-three intrinsic factors in
\eqref{geometry:intrinsicopen} are inverted.&All subsequent denominators
are units on precisely that open. Its excluded strata are not proved empty.\\
Bimodule invertibility&Both Leibniz rules and star preservation are required.&
$\sigma^{-1}=D\sigma D$ follows; no Yang--Baxter equation is added.\\\hline
\end{tabular}
\end{center}

The independent assumptions in this table define the stated restricted
problem. In particular the trace, support and regularity choices are
not consequences of being a QLC. Conversely none is dropped when
the four-variable scheme is counted. Fixing the four scales to one
selects a metric in that problem; it is not a gauge quotient.

\subsection{An explicit ledger for every inverted factor}
\label{integrity:units}

There are two kinds of units. The fixed base already inverts
$2,q_i,q_i-1,\nu,\lambda_i$ and hence every $h_i(t)$.
The connection sector additionally inverts the twenty-three factors
of $\Delta_{\rm reg}$. The inventory below refers to the latter only.
It retains even the redundant path factors.

The temporal exchange is proved by all six identities in
\eqref{geometry:chartinverse}. Thus the nine units
$(d_i,q_i-d_i,q_i-d_i^2)$ are exchanged for
$(u_i,\beta_i,\alpha_i)$, with inverse formulas in both directions.
In particular $q_i-d_i$ corresponds to the inverse of $\beta_i$
times a base unit; it need not be represented by a new polynomial
factor beyond the temporal chart.

For a triangle numerator $J$ write
\[
 W_J=\prod_i\beta_i^{b_i(J)}q_i^{c_i(J)},
 \qquad \mathcal C_1(J)=W_JJ(X_1,Y_1,m).
\]
Each $W_J$ is a unit of the temporal localization. Denote the six
active triangle numerators by
\[
 J_p=\Theta_p,\quad J_{2-p}=2E-\Theta_p,\quad
 J_r=\Theta_r,\quad J_{2-r}=2E-\Theta_r,\quad
 J_s=\Theta_s,\quad J_{2-s}=2E-\Theta_s.
\]
These are the exact adjugate polynomials of the triangle system,
evaluated on its slice. The factors actually supplied in the coefficient
files, and their pullbacks, are as follows.

\begin{center}\small
\begin{longtable}{p{.08\linewidth}p{.24\linewidth}p{.58\linewidth}}
\hline
No.&Polynomial factor&Exact connection-level interpretation\\\hline
\endfirsthead
\hline No.&Polynomial factor&Exact connection-level interpretation\\\hline
\endhead
1--3&$u_i$&$u_i=d_i(q_i-1)/(q_i-d_i)$, for $i=1,2,3$.\\
4--6&$\alpha_i$&$\alpha_i=q_i(q_i-1)(q_i-d_i^2)/(q_i-d_i)^2$.\\
7--9&$\beta_i$&$\beta_i=q_i(q_i-1)/(q_i-d_i)$.\\
10&$\overline S$&$\overline S=QS_0=QS_1$.\\
11&$\Ab$&$\Ab=q_1q_2u_1u_2 A$.\\
12&$\Bb$&$\Bb=q_2q_3u_2u_3 B$.\\
13&$q_1\alpha_2u_1^2$&$\Ab\,p_0/2$.\\
14&$q_2\alpha_1u_2^2$&$\Ab\,(2-p_0)/2$.\\
15&$q_2\alpha_3u_2^2$&$\Bb\,s_0/2$.\\
16&$q_3\alpha_2u_3^2$&$\Bb\,(2-s_0)/2$.\\
17&$\mathcal C_1(E)$&$W_E E_1$.\\
18&$\mathcal C_1(J_p)/2$&$W_{J_p}E_1 p_1/2$.\\
19&$\mathcal C_1(J_{2-p})/2$&$W_{J_{2-p}}E_1(2-p_1)/2$.\\
20&$-\mathcal C_1(J_r)/2$&$-W_{J_r}E_1 r_1/2$.\\
21&$\mathcal C_1(J_{2-r})/2$&$W_{J_{2-r}}E_1(2-r_1)/2$.\\
22&$\mathcal C_1(J_s)/2$&$W_{J_s}E_1 s_1/2$.\\
23&$-\mathcal C_1(J_{2-s})/2$&$-W_{J_{2-s}}E_1(2-s_1)/2$.\\\hline
\end{longtable}
\end{center}

For instance, the factor in row 20 differs from the intrinsic active
entry $r_1$ by the already invertible $-W_{J_r}E_1/2$.
The same statement holds for every active-entry row. The temporal
inverse formulas first establish the temporal chart, rows 10--12
establish $S,A,B$, and row 17 establishes $E_1$. Rows 13--16 and
18--23 then establish exactly the active-edge units. This order proves
the equivalence of the opens without circular use of a denominator.

Before reality has been established, the original real scheme uses
$N(\Delta_{\rm reg})$, not an unsupported notion of a complex number
being nonzero. The identity
$z^{-1}=\bar z/N(z)$ proves the equivalence of the unit conditions
over arbitrary real test algebras. After reality,
$N(\Delta_{\rm reg})=\Delta_{\rm reg}^2$, and localizing at a square
is identical to localizing at its factor: the inverse of the factor is
$\Delta_{\rm reg}(\Delta_{\rm reg}^2)^{-1}$. Thus no reducedness
assumption is concealed in the use of an open set.

For completeness, all classes of division in the reconstruction are
listed here; a determinant used to prove a unit is not an extra
localizer.

\begin{center}\small
\begin{tabular}{p{.30\linewidth}p{.28\linewidth}p{.33\linewidth}}
\hline
Operation&Divisors&Justification\\\hline
Eliminate the $24$ square coefficients&$\nu,A_i(t)$&Fixed metric units.\\
Temporal coordinate changes&$d_i,q_i-d_i,q_i,\beta_i$&Temporal unit identities above.\\
Split the rank-three block&$2,S$&Base and row 10.\\
Path completion&$A,B,Y_{0,i}$&Rows 11--12 and temporal chart.\\
Triangle completion&$Y_{1,3}E_1$&Temporal chart and row 17.\\
Exchange scalar generators&$2,S,E_1^2,AB,(AB)^2$&The displayed inverse syzygies use units already listed.\\
Adjoin graph coordinates&$q_iq_ju_iu_j$, $q_iq_j\beta_i\beta_j$&Base and temporal chart.\\
Eliminate $m$&$4z_2$, $AB$&$z_2=\alpha_2/q_2$ and rows 11--12.\\
Centre numerator binding&$u_i,\alpha_2,\beta_i$ and nonzero rationals&Exactly the factors in \eqref{centre:source1}--\eqref{centre:source3}.\\\hline
\end{tabular}
\end{center}

The later primitive-coordinate circuit introduces further written
denominators only after proving they are units in the already identified
finite coordinate algebra. Its inverse certificates and the checks of
the same twenty-three original factors are part of the centre proof;
they do not define a smaller connection sector.

\subsection{Both directions of the preservation argument}
\label{integrity:ledger}

To specify exactly what is transported, put $R_t=P_t^2-I$.
The real reconstructed coefficients satisfy the following identities
\emph{before} the equations $R_t=0$ are imposed:
\begin{equation}\label{integrity:direct-residuals}
 E_{ppr}(t)=-h_p(t)(R_t)_{pr},\qquad
 T^i{}_{pp}(t)=-\frac{h_p(t)}{h_i(t)}(R_t)_{pi},
\end{equation}
and all off-diagonal-output metric and star components vanish.
Here the temporal entries are the exact conic formulas, both blocks
share $m$, and active spatial pairs have sum two; the inactive pair
is zero. These are statements of rational identities in the metric
and temporal localization with otherwise free spatial block entries.
They are stronger than evaluating the residuals at selected solutions.
The independent direct tensor checker described in the verification
section recomputes these identities without importing the component-table
or block-completion programs.

\begin{center}\small
\begin{tabular}{p{.20\linewidth}p{.36\linewidth}p{.35\linewidth}}
\hline
Original requirement&Forward implication&Reverse verification\\\hline
Both Leibniz rules&Translation-degree preservation forces the inner
formula with no $\gamma$ and determines $\sigma$.&The recovered
$\sigma$ is a bimodule map, and the same inner formula supplies the
rules on every vertex idempotent, hence every function.\\
Torsion zero&The allowed lower indices become symmetric; all
three-distinct coefficients vanish.&Reconstruct $C^i{}_{jk}=C^i{}_{kj}$;
then $\wedge\nabla e^i=de^i=0$ and torsion is left linear.\\
Metric compatibility&Unit-leading antisymmetric components eliminate
each $w$ once; the exhaustive table yields the conics and block equations.&
Conic identities cancel all off-diagonal outputs, and
\eqref{integrity:direct-residuals} cancels all diagonal outputs.\\
Star preservation&Mixed residuals first prove temporal reality and
reciprocity; the degree-zero matrix identity then proves all reality.&
All coefficients are real; the original star residuals vanish by
\eqref{integrity:direct-residuals}, and $\sigma D\sigma D=I$
extends basis preservation to every form.\\
Support and trace&Active units force the edge sums; the trace becomes
$\tr P_t=0$ without elimination of a branch.&Path and triangle formulas
give the exact zero/unit support, both trace constraints and the same $m$.\\
Open and multiplicity&The norm localization and every unit pivot are
retained; no radical is taken.&The factor ledger and inverse formulas
give the original localization over every test algebra.\\\hline
\end{tabular}
\end{center}

The full forward chain has specific algebraic witnesses. The $24$
antisymmetric metric equations have distinct unit pivots; the temporal
relations give $dd'=1$ and force each imaginary coefficient to vanish;
the two matrix inverses give the common $m$; and the path and triangle
syzygies give inverse transformations of equation ideals. Uniqueness
of both completions and of the eliminated $w$ values then proves
that forward and reverse maps compose to the identity on every
coefficient. These are the steps proved in Theorem~\ref{geometry:bridge},
not additional assumptions on the counted points.

Equivalently, for every commutative $K$-algebra $B$ there is a natural
bijection
\[
 \mathcal X(B)\ \longleftrightarrow\
 \left\{\begin{array}{c}
 (\nabla,\sigma)\text{ satisfying the original two Leibniz rules,}\\
 \wedge\nabla=d,\quad\nabla g=0,\quad
 \nabla(\omega^*)=\sigma D\nabla\omega,\\
 \text{and precisely the homogeneity, support, trace and units above}
 \end{array}\right\}.
\]
On the left the base change from $R_{\mathbb Z}$ to $K$ is understood.
This statement applies to dual numbers and other nonreduced algebras:
it preserves infinitesimal multiplicities, not just field-valued
solutions. Consequently each of the $204$ real points later counted
reconstructs an actual pair $(\nabla,\sigma)$ satisfying the original
QLC conditions. The $512$ complex points belong to the complexification
of this real coefficient scheme; no analytic conjugation condition
is imposed independently on its arbitrary complex points.

\subsection{Why this path--triangle sector is geometrically interesting}
\label{geometry:motivation}

Define the \emph{spatial coupling graph} on the three marked directions
$\{1,2,3\}$ by the active $v_{ij}$ pairs. This is a graph of
connection coefficients, not a replacement for the sixteen-vertex
base graph of the differential calculus. Its path $P_3$ at slice zero
is minimally connected on those three directions; its triangle $K_3$
at slice one adds exactly one edge and has one independent graph cycle.
That is a controlled change from tree mixing to cyclic mixing within
one unchanged geometry. A cycle of this coefficient graph alone
proves no curvature or holonomy statement.

The choice changes the completion mechanism in a precise way. The
path has five spatial unknowns. Its left and right eigenvector equations
determine them, so trace zero becomes a remaining compatibility equation
$L=0$. The triangle has six spatial unknowns. The eigenvector equations
have a universal dependence, and the trace completes the invertible
six-by-six system with determinant $Y_{1,3}E_1$. The units $A,B,E_1,S$
specify where these mechanisms are reversible. Their complements
remain excluded strata, whose QLC classification is a separate problem.

The slices are not chosen independently: the original temporal metric
and star equations force the same $m$ and
$X_{0,i}Y_{0,i}=X_{1,i}Y_{1,i}=\alpha_i/q_i$.
Thus a path solution and a triangle solution cannot simply be paired
arbitrarily. Both supports participate in the same reconstructed
four-dimensional bimodule connection. In fact every temporal--spatial
coefficient $v_{0i}(t)$ is a unit by \eqref{geometry:temporal}.
Corollary~\ref{geometry:noproduct} consequently excludes an ordinary
coordinate-product construction, even after a permutation of the four
direction labels.

There is also an obstruction to restricting these connections to a
coordinate square. Let $\pi$ be the standard calculus projection to
that face: fix the other bits in functions, retain directions $(0,i)$
and kill the other one-forms. For a discarded spatial direction $j\ne i$,
\eqref{geometry:wreconstruction} gives the exact coefficient
\[
 \operatorname{coeff}_{e^0\otimes e^0}
       ((\pi\otimes\pi)\nabla e^j)
 =C^j{}_{00}=w_{j0}(t)=-\frac{\nu x_j(1-t)}{A_j(t)}.
\]
These coefficients depend only on $t$, so fixing the other bits leaves
this formula unchanged.
The displayed coefficient is a unit by \eqref{geometry:temporal}, so
the kernel of the quotient on forms is not preserved by the connection.
Thus there is no induced connection on that coordinate-square quotient.
At the metric-data level such a square has one constant coefficient and
one nonconstant coefficient; it lies in the intermediate metric regime
mentioned in \cite[Example~8.20, p.~598]{orig-book}. This observation
does not exclude specializing the book's generic square formulas to
subfamilies, or more general extension constructions. The quotient and
ordinary-product obstructions concern the specified coordinate maps.

\begin{proposition}[Balanced zero-translation tensor block]
\label{geometry:balanced}
On the declared real coefficient scheme,
$\Pi_{t,\pm}=(I\pm P_t)/2$ are complementary idempotents of rank two.
On the degree-zero tensor submodule generated by $e^i\otimes e^i$,
the actual bimodule map has matrix
\[
 \sigma_{t,0}=\mathscr Q_t^{\mathsf T}=-H_tP_tH_t^{-1}.
\]
It satisfies $\sigma_{t,0}^2=I$, $\tr\sigma_{t,0}=0$ and
$\det\sigma_{t,0}=1$, and its two eigensummands have rank two.
\end{proposition}
\begin{proof}
The block equations give the complementary idempotents. Their traces
are two, so their images are finite projective modules of rank two:
over a residue field an idempotent has integer rank between zero and
four, and equality of that rank with its trace two permits only rank
two when $2$ is invertible. This proves the rank statement even over
nonreduced test algebras. The component formula for $\sigma$ gives
the degree-zero matrix $\mathscr Q_t^{\mathsf T}$. Transposing
$\mathscr Q_t=-H_t^{-1}P_t^{\mathsf T}H_t$ gives the stated conjugacy.
The rank-two splitting then gives the trace and determinant assertions.
\end{proof}

This gives the trace restriction a direct spectral interpretation
in a tensor block of the connection. It does not declare $P_t$ itself
to preserve the translation degrees of individual one-forms, or claim
that either splitting is orthogonal for the metric.

The mathematical motivation is therefore the coexistence of minimal
tree-versus-cycle spatial mixing, coupled temporal data, a balanced
tensor block and reversible reconstruction of genuine QLCs. At the
specified rational metric, the theorem supplies an exact example of
multiple star-preserving QLCs for one nondegenerate metric. It provides
a controlled family for subsequent connection-dependent computations.
The worked real connection at the end of Section~\ref{centre:section}
provides a concrete instance of these coupled blocks and their active entries.
No curvature calculation, action principle, continuum limit or
physical uniqueness principle selecting this particular support has
been established by the count.

Finally, marking must not be confused with a multiplicative symmetry
argument. The permutation group $S_3$ produces three distinct
\emph{unoriented} path supports, since swapping its two endpoints
stabilizes a path. There are six mixed support patterns only if the
choice of which slice carries the path is included. This is a
combinatorial statement, not a new count of QLCs. A permutation of
spatial directions usually also permutes the anisotropic metric
parameters, so it need not act on the same fixed-metric fibre. No
factor of six, quotient count or enlarged-sector real count is
inferred here.

\section{The arithmetic control fibre and the generic degree}
\label{generic:section}

The generic degree will be proved from one finite arithmetic fibre, but a
count of solutions in that fibre alone is insufficient. We must exclude
the entire boundary of the \emph{same family}, including limits with
varying metric parameters, and check its relative Jacobian. This section
gives those two additional arguments. All finite-field computations below
are in characteristic $p_*=65537$, not characteristic two.

\subsection{A common arithmetic presentation and its proper closure}
\label{generic:model}
Set
\[
 R=\mathbb Z[q_1,q_2,q_3,
 (2q_1q_2q_3(q_1-1)(q_2-1)(q_3-1))^{-1}],\qquad
 b_*=(p_*,q_1-4,q_2-9,q_3-49)\in\Spec R.
\]
The inverted polynomial has residue $962\ne0$. In this section the
control metric is $q_*=(4,9,49)$ over $\mathbb F_{p_*}$; it is distinct
from the rational characteristic-zero metric $q^0$ used later.
Write
\[
 \alpha_i=q_i-(u_i-1)^2,\quad \beta_i=q_i+u_i-1,\quad Q=q_1q_2q_3,
 \quad D_\beta=Q\beta_1\beta_2\beta_3,
\]
\[
 \Ab=q_2\alpha_1u_2^2+q_1\alpha_2u_1^2,\qquad
 \Bb=q_3\alpha_2u_3^2+q_2\alpha_3u_2^2,
\]
and put
\[
 D_A=q_1q_2u_1u_2,\quad D_B=q_2q_3u_2u_3,\quad
 D_{ij}=q_iq_j\beta_i\beta_j,
 \quad N_{ij}=q_i\alpha_i\beta_j^2+q_j\alpha_j\beta_i^2.
\]
The five graph equations are
\[
 E_A=D_A\xi-\Ab,\quad E_B=D_B\eta-\Bb,\quad
 E_{ij}=D_{ij}c_{ij}-N_{ij}\quad(ij=12,13,23).
\]
The remaining four graph equations are
\begin{align}
 H_p&=q_2\xi\eta+4\alpha_2=0,\label{generic:Hp}\\
 L_p&=Q\xi\eta m+\xi R_{23}+\eta R_{12}=0,\label{generic:Lp}\\
 T&=Q(m^2-1)+q_2q_3\alpha_1+q_1q_3\alpha_2+q_1q_2\alpha_3=0,
 \label{generic:top}\\
 H_{\rm graph}&=-Q^4\bigl(H_0(z,c)+mH_1(z,c)\bigr)=0.\label{generic:trianglegraph}
\end{align}
Here $z_i=\alpha_i/q_i$, and
\[
 R_{23}=q_1q_3\alpha_2+q_1q_2\alpha_3,\qquad
 R_{12}=q_2q_3\alpha_1+q_1q_3\alpha_2.
\]
Here $H_0,H_1$ are the explicit invariant polynomials of the preceding
reconstruction, not additional unspecified equations. Thus all nine
equations are integer polynomials. Denote their ideal by $I$.

Use the full 23-factor localization product $s$ from the declared sector.
Its factors are the nine $u_i,\alpha_i,\beta_i$; the three
$\Ab,\Bb,Q\sum_i\alpha_i/q_i$; the four path-entry numerators; and
the seven triangle factors consisting of the cleared completion
determinant and six active-entry numerators. Each is the polynomial
obtained in the source reconstruction, including the stated unit
normalizations. Let
\[
 G=\Spec\bigl((R[u_1,u_2,u_3,m,\xi,\eta,c_{12},c_{13},c_{23}]/I)_s\bigr).
\]
No factor is dropped in the argument.

Put $h=u_1u_2u_3\beta_1\beta_2\beta_3$. On $D(h)$ all five graph
denominators are units. Eliminating their unique sections identifies $G$
scheme-theoretically with a localization $X$ of a four-equation family
in $(u_1,u_2,u_3,m)$. Useful unnormalized source equations are
\begin{align}
 P_H&=\Ab\Bb+4Q\alpha_2u_1u_2^2u_3,\label{generic:PH}\\
 P_L&=\Ab\Bb m+
 u_2u_3(q_3\alpha_2+q_2\alpha_3)\Ab
 +u_1u_2(q_2\alpha_1+q_1\alpha_2)\Bb,\label{generic:PL}\\
 P_{\rm top}&=T,\label{generic:PT}\\
 \Phi&=D_\beta^4\,
  \mathcal H\left(\frac{\alpha_1}{\beta_1},
                    \frac{\alpha_2}{\beta_2},
                    \frac{\alpha_3}{\beta_3},
                    \frac{\beta_1}{q_1},
                    \frac{\beta_2}{q_2},
                    \frac{\beta_3}{q_3},m\right),\label{generic:Phi}
\end{align}
where $\mathcal H(X,Y,m)=H_0(X_iY_i,X_iY_j+X_jY_i)
+mH_1(X_iY_i,X_iY_j+X_jY_i)$. Each displayed clearing multiplier is a
unit on $X$. In particular, the support is not replaced by the support
of a larger unsaturated compact system.

Take the affine scheme closure of $X$ by saturating its four-equation
ideal by the full localization product over $R$, and let
\[
 \overline X\subset\mathbb P^4_R,
 \qquad[W:U_1:U_2:U_3:M],
\]
be its projective scheme closure. Both closures are formed \emph{before}
specialization to $b_*$. The morphism $\overline X\to\Spec R$ is
projective and finitely presented, hence proper. We prove
\begin{equation}\label{generic:controlgoal}
 \overline X_{b_*}=X_{b_*},\qquad
 X_{b_*}\text{ is geometrically reduced of length }512,
\end{equation}
and verify the actual relative Jacobian on this fibre.

\subsection{Exact arithmetic certificates and their meaning}
\label{generic:certificates}
The supplement under \texttt{certificates/generic/} supplies the full
polynomials, not numerical samples of their values. A sparse polynomial
is a list of integer coefficient--exponent pairs with an explicit
variable order; univariate residues are in ascending coefficient order.
Arithmetic in $\mathbb F_{p_*}[t]/(E)$ always means polynomial reduction
modulo the displayed $E$. It is not assumed that $E$ is irreducible.

The verifier \texttt{audit/generic\_certificate.py} reconstructs the
equations from \eqref{generic:Hp}--\eqref{generic:Phi}, computes the
stated substitutions, divisions and determinants, and checks every
identity below. In particular its acceptance is based on equalities of
polynomial coefficients and explicit unit identities, not on stored
truth values or file digests. The finite data used in the proof have
the following roles:
\begin{center}\small
\begin{tabular}{@{}lp{0.67\textwidth}@{}}\toprule
Certificate & Exact mathematical content\\\midrule
Finite chart boundary & All 26 nontrivial $u/\beta$ support words;
four mixed B\'ezout identities, six single-$\beta$ B\'ezout identities,
three double-$\beta$ unit ideals, and two endpoint quotient units.\\
Exceptional lines & Eight complete double-$\alpha$ lines; the global
saturation identity and one univariate B\'ezout identity per line.\\
Finite projection & Exact pseudo-remainder and resultant identities;
the squarefree projection polynomial of degree 642.\\
Projection exclusion & Every factor of its degree-130 complement,
including all degree-drop and zero-slope branches.\\
Finite rational representation & A squarefree degree-512 polynomial,
all coordinate residues, original equation residues and unit tests.\\
Proper boundary & All global source coefficients, continuous pole-cone
faces, two global colon identities and six signed one-pole eliminations.\\
Graph Jacobian & All nine original graph equations, all 23 original
factors, full $9\times9$ and $4\times4$ determinants and their inverses.\\
\bottomrule\end{tabular}
\end{center}
Finite polynomial arrays too long to print are part of these certificates.
The verification rules and the exhaustiveness argument are given next;
the arrays are witnesses for those rules, not replacements for them.

\subsection{All finite graph-chart boundary supports}
\label{generic:chartboundary}
Let $\Gamma$ be the closure of the graph of $X$ in
$\overline X\times_R(\mathbb P^1_R)^5$. Its projection to
$\overline X$ is proper and surjective on supports: its closed image
contains the dense open $X$. Thus every proposed finite source boundary
point has a graph lift after a residue-field extension.

For a point in the closure of a finite-type open immersion over a
Noetherian base, the DVR specialization lemma supplies a trait whose
generic point lies in the open and whose closed point is the proposed
point. We use this with arbitrary residue extensions of
$\mathbb F_{p_*}$. The parameters of the trait reduce to $q_*$ but may
vary along it. All the nonzero residue constants used below are therefore
units on the trait. This standard closure lemma is recorded in
\href{https://stacks.math.columbia.edu/tag/0CM2}{Stacks, Tag 0CM2}.

At a finite source point $u_i,m$ are integral. At the control point
$\beta_i-u_i$ has residue $3,8,48$ for $i=1,2,3$ respectively. Hence
$u_i$ and $\beta_i$ cannot both vanish at one index. Use the letters
$U,B,N$ for $u_i=0$, $\beta_i=0$, and neither zero. The complement of
$D(h)$ consists of exactly the $3^3-1=26$ words other than $NNN$.

First consider a vanishing endpoint $u_1$ with $u_2,u_3$ units. If
$a=v(u_1)>0$, then $\Ab$ is a unit, $v(\xi)=-a$, and
\eqref{generic:Hp} gives
$v(\eta)=v(\alpha_2)+a\ge a$. The equation $E_B=0$ forces
$\Bb=0$ in the residue field. Dividing \eqref{generic:Lp} by $\xi$
on the generic point gives
\[
 R_{23}=-Q\eta m-(\eta/\xi)R_{12}.
\]
Both terms on the right have positive valuation, so $R_{23}=0$ in the
residue field as well. The reflected endpoint gives
$\Ab=R_{12}=0$. These statements allow arbitrarily unequal positive
orders of the vanishing factors.

The seven words with $U$ letters and no $B$ letters are excluded as follows:
\begin{center}\small
\begin{tabular}{@{}ll@{}}\toprule
Word & Contradiction\\\midrule
$UNN$ & $\Bb=R_{23}=0$ and the left endpoint quotient unit below.\\
$NNU$ & $\Ab=R_{12}=0$ and the right endpoint quotient unit below.\\
$NUN$ & Both path coordinates have poles, contradicting $H_p=0$.\\
$UNU$ & Both path coordinates have poles, contradicting $H_p=0$.\\
$UUN$ & In $L_p$, $\eta R_{12}$ is the unique pole;
 $R_{12}(0,0,u_3)=2891\ne0$.\\
$NUU$ & The reflected unique-pole argument;
 $R_{23}(u_1,0,0)=3296\ne0$.\\
$UUU$ & The unique finite triangle section contradicts the two trace equations.\\
\bottomrule\end{tabular}
\end{center}
For the two endpoint quotient calculations, form
\[
 A_L=\mathbb F_{p_*}[u_2,u_3]/(\Bb,R_{23})|_{u_1=0},\qquad
 A_R=\mathbb F_{p_*}[u_1,u_2]/(\Ab,R_{12})|_{u_3=0}.
\]
Both have dimension eight. On the $\beta$-open chart the triangle sections
are regular; substitute them into $H_{\rm graph}$ and eliminate $m$
against $T$. After clearing only unit denominators, the obstruction classes
$g_L,g_R$ have invertible multiplication maps with determinants
\[
 \det M_{g_L}=35553,\qquad \det M_{g_R}=61639
 \quad\text{in }\mathbb F_{p_*}.
\]
The certificate supplies the two eight-element monomial bases, all reduced
coefficients and inverse classes satisfying $g_Lg_L^{-1}=g_Rg_R^{-1}=1$.
A class with an inverse cannot vanish at a geometric point. This argument
does not assume the quotient algebras are reduced.
At $UUU$, the triangle sections are
\[
 (c_{12},c_{13},c_{23})=(59/36,339/196,824/441).
\]
After nonzero scalar normalizations the two trace equations are
$4924-13357m=0$ and $1764m^2+2855=0$. The first forces $m=1070$,
and the second has residue $18263\ne0$ there.

There are twelve mixed $U/B$ words. The words $BUN,NUB$ give the middle
double-pole contradiction; $UUB,UBU,BUU$ are the adjacent or separated
$U$ contradictions just proved. At $UBB$ and $BBU$, respectively one
path coordinate has a pole and the other has the finite nonzero residue
$-56$ and $-11$; at $BUB$ both have poles. These contradict $H_p$.
The four remaining words reduce to endpoint necessary pairs:
\begin{center}\small
\begin{tabular}{@{}llr@{}}\toprule
Word & Necessary univariate pair & Resultant modulo $p_*$\\\midrule
$UBN$ & $(\Bb,R_{23})|_{u_1=0,u_2=-8}$ &49478\\
$UNB$ & $(\Bb,R_{23})|_{u_1=0,u_3=-48}$ &62173\\
$BNU$ & $(\Ab,R_{12})|_{u_3=0,u_1=-3}$ &33484\\
$NBU$ & $(\Ab,R_{12})|_{u_3=0,u_2=-8}$ &13379\\
\bottomrule\end{tabular}
\end{center}
For each pair $(f,g)$, the supplied multipliers satisfy $af+bg=1$.
For example the $UBN$ pair, with indeterminate $v=u_3$, is
\[
 f=-4104v^2+1152v+27648,\qquad
 g=-36v^2+72v-12384,
\]
and its identity is
\[
 (30332v+32123)f+(15613v+26845)g=1
 \quad\text{in }\mathbb F_{p_*}[v].
\]
All four identities are verified coefficientwise and therefore exclude
points over every extension field.

Next suppose that at least two $\beta_i$ vanish and no $u_i$ vanishes.
For each of the three pairs, substitute $u_i=1-q_i$ into the five path
equations $E_A,E_B,H_p,L_p,T$. The resulting polynomials generate the
unit ideal in the remaining $u$ coordinate, $m,\xi,\eta$ over
$\mathbb F_{p_*}$. The certificate lists these five polynomials and the
exact finite-field elimination; its reduced Gr\"obner basis is $[1]$ in
each case. The triple-$\beta$ word lies in these pair loci, so all four
multiple-$B$ words are covered. No triangle coordinate is assumed finite
in this calculation.

It remains to consider the three words with exactly one $B$ and no $U$.
If $\beta_i$ vanishes, its two incident triangle sections have poles,
and the third is integral. From the graph equations,
\[
 \beta_i c_{ij}=\frac{\alpha_i\beta_j}{q_j}
 +\frac{\alpha_j\beta_i^2}{q_i\beta_j}.
\]
Here $\alpha_i$ has the nonzero residue $q_i(1-q_i)$. The degree-four
part of the triangle equation in the incident pair, up to a nonzero base
unit, is $314(c_{ij}^2-c_{ik}^2)^2$. Multiplying the global equation by
$\beta_i^4$ and taking its residue therefore forces
\[
 (\beta_j/q_j)^2=(\beta_k/q_k)^2,
 \qquad\beta_j/q_j=\varepsilon\beta_k/q_k,
 \quad\varepsilon\in\{1,-1\}.
\]
Thus there are precisely six necessary signed branches. Substitute each
linear relation into the path equation and the two trace equations and
clear the already unit coordinate denominators. The resulting path
polynomial $f$ and division-free trace resultant $g$ have degrees
\begin{center}\small
\begin{tabular}{@{}lccc@{}}\toprule
Vanishing factor & $\beta_1$ & $\beta_2$ & $\beta_3$\\\midrule
Degrees $(\deg f,\deg g)$, for each sign &(6,14)&(4,10)&(6,14)\\
\bottomrule\end{tabular}
\end{center}
Each of the six pairs has a supplied identity $af+bg=1$. In forming $g$
no leading coefficient of the linear trace equation is inverted: if the
two equations are $am+b=0$ and $cm^2+d=0$, the identity
\begin{equation}\label{generic:traceelimination}
 a^2(cm^2+d)-c(am-b)(am+b)=cb^2+a^2d
\end{equation}
proves the necessary equation even where $a=0$. The certificate reconstructs
these six pairs from the original substituted equations before verifying
their B\'ezout identities. All $7+12+4+3=26$ words are now excluded.
Consequently the finite-source graph closure has no point on $V(h)$.

\subsection{The eight exceptional lines on the remaining finite chart}
\label{generic:exceptionallines}
On $D(h)$ the graph coordinates are uniquely reconstructed as schemes.
However, one must not divide a resultant factor that can vanish there.
At the control point write $a_i=\alpha_i$. On $D(h)\cap V(P_H)$,
\begin{equation}\label{generic:exceptionalsplit}
 V(a_2\Bb)=V(a_1,a_2)\ \cup\ V(a_2,a_3).
\end{equation}
Indeed, $\Bb=0$ in \eqref{generic:PH} forces $a_2=0$, and then
$P_H=q_2^2u_2^4a_1a_3$. Conversely either pair makes $P_H$ vanish.
Thus the exceptional set is exactly the four lines
\[
 (u_1,u_2)\in\{-1,3\}\times\{-2,4\}
\]
and the four lines
\[
 (u_2,u_3)\in\{-2,4\}\times\{-6,8\}.
\]
Their intersections are retained. On the complement of
\eqref{generic:exceptionalsplit}, $a_2,\Ab,\Bb$ are all units.

To exclude each \emph{whole} line in the global closure we need an
equation in the global saturated ideal, not saturation after specialization.
Define
\[
 C=4u_2^2m+4u_2(u_1+u_3)
    +\xi(u_3^2-u_2^2)+\eta(u_1^2-u_2^2).
\]
Direct expansion over the integer polynomial ring gives
\begin{align}
 Q\xi\eta C-4u_2^2L_p
  ={}&4q_3\eta E_A+4q_1\xi E_B\notag\\
   &+q_1q_3[\eta(u_1^2-u_2^2)+\xi(u_3^2-u_2^2)]H_p=:W_0,
   \label{generic:coreidentity}\\
 D_B\eta E_A+\Ab E_B
  ={}&D_AD_B\xi\eta-\Ab\Bb.\label{generic:sectionidentity}
\end{align}
Combining them yields
\[
 Q\Ab\Bb C
 =D_AD_B(4u_2^2L_p+W_0)
   -QC(D_B\eta E_A+\Ab E_B)\in I.
\]
Since $Q,\Ab,\Bb$ are inverted on $G$, this proves $C\in I:s^\infty$
over $R$. Every point of the specialized global closure must satisfy its
specialization. No saturation--specialization interchange has been used.

On each of the eight lines, $4u_2^2$ is a unit. Substitute the five graph
sections into $C=0$ and solve uniquely for $m=M(v)$, with $v$ the free
coordinate. Substitute this expression into $T$ and the actual triangle
equation. Cancellation leaves denominators consisting of nonzero constants
and powers of $v$, which are units on $D(h)$. The numerator pairs have
the following degrees:
\begin{center}\small
\begin{tabular}{@{}lcrr@{}}\toprule
Fixed coordinates & Free coordinate & Top numerator & Triangle numerator\\\midrule
$(u_1,u_2)=(-1,-2)$&$u_3$&4&9\\
$(u_1,u_2)=(-1,4)$&$u_3$&4&9\\
$(u_1,u_2)=(3,-2)$&$u_3$&4&9\\
$(u_1,u_2)=(3,4)$&$u_3$&4&9\\
$(u_2,u_3)=(-2,-6)$&$u_1$&3&9\\
$(u_2,u_3)=(-2,8)$&$u_1$&4&9\\
$(u_2,u_3)=(4,-6)$&$u_1$&4&9\\
$(u_2,u_3)=(4,8)$&$u_1$&3&9\\
\bottomrule\end{tabular}
\end{center}
For example the first line gives
$M(v)=(199v^2-202v-144)/(392v)$. The complete numerator polynomials
and multipliers $A(v),B(v)$ are supplied for all eight lines; the verifier
reconstructs them from \eqref{generic:coreidentity} and checks
\[
 A(v)T_C(v)+B(v)H_C(v)=1
 \quad\text{in }\mathbb F_{p_*}[v].
\]
Thus every exceptional line, including every intersection, is empty in
the specialized global closure.

\subsection{An exhaustive finite projection on the allowed complement}
\label{generic:projection}
For this subsection only, let $P_H^c,P_0^c,P_\triangle^c$ denote the
three primitive compact control polynomials supplied in
\texttt{certificates/generic/compact/}. They are defined from the four
source equations by eliminating the scalar $m$ without division. More
explicitly, write the normalized source equations as
\[
 L=A_m m+B_m,\qquad \Phi_0=C_m m+D_m,
 \qquad T_0=T_2m^2+T_{00}.
\]
The defining polynomial identities are
\begin{align}
 P_H&=2P_H^c,\notag\\
 A_m\Phi_0-C_mL&=A_mD_m-C_mB_m=2P_\triangle^c,
 \label{generic:compacttriangle}\\
 A_m^2T_0-T_2(A_mm-B_m)L
   &=T_2B_m^2+T_{00}A_m^2=-4P_0^c.\label{generic:compacttop}
\end{align}
For these compact identities the normalized polynomials are
$L=P_L|_{q_*}$, $T_0=-P_{\rm top}|_{q_*}$, and $\Phi_0=\Phi|_{q_*}$;
the path polynomial is $P_H^c=P_H|_{q_*}/2$. The nine-variable graph
triangle equation, after substituting its sections and multiplying by
$D_\beta^4$, equals $-Q^4\Phi$; here $Q=1764$ and
$-Q^4=-9682651996416\ne0\pmod{p_*}$. The verifier checks all these
normalizations against the displayed graph equations, not merely against
the compact data.
They give necessary compact equations even where $A_m=0$.

Work on $D(h\alpha_2\Bb)$. In the variable $u_1$, the exact resultant is
\begin{equation}\label{generic:firstresultant}
 \operatorname{Res}_{u_1}(P_H^c,P_0^c)
 =13108\,u_2^8\alpha_2^4\Bb^2 F(u_2,u_3)
 \quad\text{in }\mathbb F_{p_*}[u_2,u_3].
\end{equation}
Only the displayed units are cancelled. The last nonconstant
subresultant is $\ell_1u_1+\ell_0$; coefficient comparison also proves
that it is the negative pseudo-remainder of $P_0^c$ by $P_H^c$.
Thus every common compact root satisfies it, including specialized
leading-coefficient degeneracies. Define the polynomial
\[
 C_{\rm proj}=\sum_{i=0}^{10}[u_1^i]P_\triangle^c
                       (-\ell_0)^i\ell_1^{10-i}.
\]
If $\ell_1\ne0$, this is the denominator-cleared substitution
$\ell_1^{10}P_\triangle^c(-\ell_0/\ell_1,u_2,u_3)$.
If $\ell_1=0$, the necessary linear relation gives $\ell_0=0$;
the same polynomial still vanishes. Hence every allowed compact point
projects to $F=C_{\rm proj}=0$.

The exact resultant of these two polynomials has $u_3$-degree 8568.
Its monic squarefree part factors as
\begin{equation}\label{generic:projectionfactor}
 \operatorname{rad}\operatorname{Res}_{u_2}(F,C_{\rm proj})
   =E(u_3)J(u_3),\qquad \deg E=512,\quad\deg J=130,
 \quad\gcd(E,J)=1.
\end{equation}
This is a coefficient identity over $\mathbb F_{p_*}$, and $E,J$ are
both supplied coefficientwise. The equality is used as a necessary
projection cover, not as an assertion that all its roots lift.

Here is the complete exclusion rule for the complementary factor $J$.
Factor $J$ into its 37 irreducible factors $j$. For each, work in the
finite field $k=\mathbb F_{p_*}[z]/(j)$ and compute in $k[y]$
\[
 G_y=\operatorname{rad}\gcd(F,C_{\rm proj}),\qquad
 D_y=\gcd(G_y,\ell_1),\qquad L_y=G_y/D_y.
\]
These split every projected root into the zero-slope and nonzero-slope
charts. For linear $j$, factor the resulting polynomials over
$\mathbb F_{p_*}$; for the other factors the resulting nonconstant
polynomials in $y$ are linear, as checked by the certificate. In each
resulting residue field, substitute $y,z$ into \emph{all three} compact
polynomials and form the radical common polynomial in $u_1$.
If they all vanish identically, the certificate checks that one of
$u_2,u_3,\alpha_2,\beta_2,\beta_3$ vanishes identically, so the branch
is outside the allowed complement. Otherwise remove its common roots
with
\[
 u_1,u_2,u_3,\alpha_2,\beta_1,\beta_2,\beta_3
\]
by repeated polynomial gcd and exact division. The remainder is a nonzero
constant in every case. Thus there is no allowed $J$-root. The weighted
degrees of the zero-slope and nonzero-slope projection parts are 129 and
4 respectively. They count projected $(y,z)$ branches, so several $y$
roots can lie over one root of $J$; they need not sum to $\deg J$ and
are not counts of admissible source points.
The certificate gives the full 37-factor inventory and every resulting
branch, so neither a vanished slope nor an identically zero compact
specialization is silently discarded.

For the $E$ part, form the generally non-field algebra
\[
 A_*=\mathbb F_{p_*}[t]/(E(t)).
\]
All Euclidean elimination pivots used below are certified units in $A_*$.
The monic common divisor of $F(y,t),C_{\rm proj}(y,t)$ is linear in $y$,
giving a unique $u_2=r_2(t)$. The monic common divisor of
$P_H^c(u_1,r_2,t),P_0^c(u_1,r_2,t)$ is linear in $u_1$, giving a unique
$u_1=r_1(t)$. Finally $A_m(r_1,r_2,t)$ is a unit and gives
$m=r_m(t)=-B_m/A_m$. The three degree-below-512 residue arrays are
included in the certificate. Direct substitution proves every compact
and original source equation. Moreover,
\[
 \deg E=512,\quad\gcd(E,E')=1,
\]
and all 23 original localization factors evaluate to units in $A_*$.
For each factor $f$, this is checked as a polynomial inverse identity
$f(t)v_f(t)\equiv1\pmod E$, not as a nonzero value at one sample.
The source Jacobian is a unit as well, as detailed below.

The two linear gcd identities prove uniqueness of every allowed lift
above each root of $E$. Conversely the explicit residues provide such
a lift for every root. Thus exactly 512 geometric points remain. The
eight exceptional lines and all 26 chart supports have already been
excluded. Since $I:s^\infty$ and $I$ agree after localization at $s$,
the reconstructed points lie on the specialized global closure with the
same local scheme structure. The Jacobian unit proves that these local
rings are reduced, not merely that their support has size 512.

\subsection{Genuine source infinity, including moving-base limits}
\label{generic:infinity}
Suppose a proposed point of $\overline X_{b_*}$ has $W=0$. Apply the
DVR closure lemma as above. At least one of $u_1,u_2,u_3,m$ has a pole:
otherwise $[1:u_1:u_2:u_3:m]$ extends in the affine chart, and
separatedness of projective space forbids a different projective limit.
Set
\[
 r_i=\max(0,-v(u_i)),\qquad r_m=\max(0,-v(m)).
\]
The four global polynomials \eqref{generic:PH}--\eqref{generic:Phi}
have respectively 23,55,8,1198 monomials after grouping by exponents
in $(u_1,u_2,u_3,m)$. Every grouped coefficient specializes to a nonzero
element at $b_*$. This is checked for every coefficient, including all
164616 integer monomials in the global expansion of $\Phi$.
Consequently these coefficients are units on any moving-base trait;
the following initial-face arguments do not require the parameters to
be constant along that trait. We express the control-field faces in
the scalar normalization $P_H^c=P_H|_{q_*}/2$,
$T_0=-P_{\rm top}|_{q_*}$ and $\Phi_0=\Phi|_{q_*}$; all these changes are
by nonzero control-field scalars. The roles and normalizations are:
\begin{center}\small
\begin{tabular}{@{}lll@{}}\toprule
Role & Control polynomial / checker key & Centre numerator / file\\\midrule
Path & $P_H|_{q_*}/2=P_H^c$ / \texttt{PH4} & $P_p$ / \texttt{Hp.txt}\\
Scalar path & $P_L|_{q_*}=L$ / \texttt{PL4} & eliminated by $m=m_0$\\
Top & $-P_{\rm top}|_{q_*}=T_0$ / \texttt{FT4} & $P_F$ / \texttt{F.txt}\\
Triangle & $\Phi|_{q_*}=\Phi_0$ / \texttt{PTH4} & $P_T$ / \texttt{Htriangle.txt}\\
\bottomrule\end{tabular}
\end{center}
The control column is over $\mathbb F_{p_*}$ at $q_*$; the centre column
is over $\mathbb Q$ at $q^0$ after scalar elimination. It records the
roles of different polynomial arrays. The exact centre multipliers are
\eqref{centre:source1}--\eqref{centre:source3}. In particular the generic
top equation is denoted $P_{\rm top}$, and $P_T$ denotes the centre
triangle numerator.

The leading part of the top equation, up to a nonzero scalar, is
\[
 441u_1^2+196u_2^2+36u_3^2-1764m^2.
\]
The largest positive pole order in this equation must occur at least
twice. In particular $r_m\le\max_i r_i$, and $m$ cannot be the only
pole. If all three $u_i$ have positive pole order, the path equation has
the unique dominant monomial $377u_1^2u_2^4u_3^2$, a contradiction.

For exactly two positive $u$-pole orders, the following componentwise
maximal faces apply throughout the \emph{entire} positive rate cone:
\begin{center}\small
\begin{tabular}{@{}lll@{}}\toprule
Pole coordinates & Path factor in the finite coordinate & Triangle factor\\\midrule
$u_2,u_3$ &$13u_1^2-18u_1-27$&$(u_1+3)^4$\\
$u_1,u_3$ &$(13u_2^2-8u_2-32)(29u_2^2-49u_2-196)$&$(u_2+8)^4$\\
$u_1,u_2$ &$29u_3^2-9u_3-216$&$(u_3+48)^4$\\
\bottomrule\end{tabular}
\end{center}
Only nonzero scalar factors and powers of the pole coordinates have
been suppressed. Every other $m$-free monomial has coordinatewise exponents no larger
than the displayed maximal pair and at least one strict inequality.
Thus the same face applies throughout the positive rate cone.
The triangle face forces the finite residue to be
$-3,-8,-48$ respectively, where the displayed path factors are
$144$, $864\cdot2052$, and $67032$, all nonzero modulo $p_*$.
There are 501 $m$-linear triangle monomials. Their exponents in either
positive coordinate are at most seven, while the $m$-free maximal
face has exponents $(8,8)$. Their weighted deficit is at least
\[
 r_i+r_j-r_m\ge\min(r_i,r_j)>0.
\]
Thus arbitrary unequal rates, the equality wall, and an additional
$m$ pole cannot change this face. Finite coordinates with positive
valuation, and hence zero residue, are included by the displayed
polynomial conditions.

Finally suppose exactly one $u_i$ has positive pole order. The top
equation forces $r_m=r_i$ and, at the leading coefficient level,
\[
 m/u_1=\pm1/2,\qquad m/u_2=\pm1/3,\qquad m/u_3=\pm1/7.
\]
For any global polynomial
$f=\sum_{\alpha,\ell}a_{\alpha,\ell}(q)u^\alpha m^\ell$ occurring here,
define the one-pole chart polynomial by
\[
 d_i(f)=\max_{a_{\alpha,\ell}\ne0}(\alpha_i+\ell),\qquad
 f^{\infty,i}(\widehat u_i,M)=
 \left.t^{d_i(f)}f(u_i=t^{-1},m=M/t,\widehat u_i;q)
 \right|_{t=0,\,q=q_*}.
\]
Here $\widehat u_i$ denotes the two remaining $u$ coordinates. All grouped
coefficients are units on the moving-base trait, so the same weighted
support applies before and after reduction. This formula is polynomial
in the finite coordinates, including zero residues, and defines
$\Phi^{\infty,i}$ in particular.

For an endpoint $i=1$ or $3$, put $t=1/u_i$ and $M=m/u_i$.
The other source coordinates and $M$ are integral. Let
\[
 F_H=t^2P_H,\qquad F_L=t^3P_L,\qquad
 E_1=\Bb,\quad E_3=\Ab,
\]
\[
 D_i=q_i\alpha_2-q_2u_2^2,\qquad N_i=MD_i-q_2u_2.
\]
At $t=0$ the two chart equations give $E_iD_i=E_iN_i=0$.
If $E_i\ne0$, then $D_i=N_i=0$ forces $u_2=0$ and then
$D_i=q_i(q_2-1)\ne0$. Therefore $E_i=0$.

Exact division over $\mathbb Z[q_1,q_2,q_3]$ gives a 269-term
polynomial $K_i$ for each endpoint satisfying
\begin{equation}\label{generic:endpointcolon}
 t u_2K_i=D_iF_L-N_iF_H.
\end{equation}
On the trait's generic point $t$ and $u_2$ are nonzero, so $K_i=0$
there and then in the valuation ring. This uses neither division by the
residue of $u_2$ nor saturation after specializing $t=0$.
For each sign of $M$, eliminate one finite coordinate from
\[
 E_i=0,\quad K_i|_{t=0}=0,\quad \Phi^{\infty,i}=0.
\]
The exact pairs of equations and the eliminated coordinate are:
\begin{center}\small
\begin{tabular}{@{}clll@{}}\toprule
$i$ & Two pairs & Eliminate & Retain\\\midrule
$1$ & $(E_1,K_1|_{t=0}),\ (E_1,\Phi^{\infty,1})$ & $u_2$ & $u_3$\\
$2$ & $(P_H^{\infty,2}/2,P_L^{\infty,2}),\ (P_H^{\infty,2}/2,\Phi^{\infty,2})$
 & $u_1$ & $u_3$\\
$3$ & $(E_3,K_3|_{t=0}),\ (E_3,\Phi^{\infty,3})$ & $u_1$ & $u_2$\\
\bottomrule\end{tabular}
\end{center}
All entries are reduced at $q_*$ and at each signed value of $M$ above;
the factor $1/2$ is the control path normalization. The endpoint
resultant degrees are 18 and 32, and the middle-chart degrees are 8 and 32.
For all six signed cases the
certificate supplies their complete coefficient arrays $A,B$ and
multipliers $U,V$ satisfying
\[
 UA+VB=1\quad\text{in }\mathbb F_{p_*}[z].
\]
The verifier first reconstructs each resultant from the actual chart
polynomials by finite-field Sylvester determinants and interpolation
under a checked degree bound, and only then checks the unit identity.
A common chart root makes both resultants vanish even if a specialized
leading coefficient drops: this is the polynomial Sylvester identity,
not division by that leading coefficient. Each signed chart is therefore
empty over the algebraic closure.

The cases of zero, one, two and three positive $u$-pole orders exhaust
every possible pole vector. We have proved
\begin{equation}\label{generic:noinfinity}
 \overline X_{b_*}\cap V(W)=\varnothing.
\end{equation}

\subsection{The actual relative Jacobian and the finite \texorpdfstring{\'etale}{etale} neighbourhood}
\label{generic:jacobian}
In the algebra $A_*=\mathbb F_{p_*}[t]/E$, reconstruct all nine
coordinates by the preceding rational representation and the five graph
sections. Differentiate \emph{all nine original graph equations} with
respect to their nine variables, substitute these residues and compute
the determinant by a division-free subset recurrence. Explicitly, for a
matrix $J$ over $A_*$, set $d_\varnothing=1$ and, for a column subset
$S$ of size $k$,
\[
 d_S=\sum_{j\in S}(-1)^{\#\{\ell\in S:\ell>j\}}
                  J_{k,j}d_{S\setminus\{j\}}.
\]
Then $d_{\{1,\ldots,9\}}=\det J$. This uses no pivot divisions and
therefore remains valid in the generally reducible algebra $A_*$.
The resulting determinant has a supplied polynomial inverse modulo $E$.
Independently, differentiate the four actual source equations in their
four variables. Their determinant is also a unit, and elimination of the
five graph-section rows verifies the Schur-complement equality with the
nine-variable determinant. The coefficients of every one of the 23
original factors are evaluated afresh in the same algebra and each
product with its supplied inverse reduces to one. These tests certify
the actual original open and Jacobian, rather than those of a convenient
replacement system.

\begin{proposition}\label{generic:controlfibre}
The whole projective control fibre is equal to the original localized
source fibre. It is geometrically reduced of length 512, and the original
source morphism is \'etale at each of its points.
\end{proposition}
\begin{proof}
The proper graph projection reduces every finite source boundary point
to the graph boundary. Section~\ref{generic:chartboundary} excludes
$V(h)$. On $D(h)$ the graph is a regular closed section; the eight-line
argument and the exhaustive projection calculation then exclude all
remaining localization boundaries and leave exactly the reconstructed
512 admissible points. Equation~\eqref{generic:noinfinity} excludes
every point outside the finite source chart. Hence the entire support
of $\overline X_{b_*}$ lies in $X$.

An open subscheme containing every point of a scheme is the entire
scheme. Restriction of global scheme closure to an open commutes with
the closure construction, and localization commutes with base change.
Since $(I:s^\infty)_s=I_s$ before specialization, the control fibre
has precisely the source scheme structure at these points; no nilpotent
or embedded part is dropped. The source Jacobian units prove that all
these local rings are reduced. They also prove relative \'etaleness
by the square Jacobian criterion: there are nine equations in nine
graph variables, or equivalently four equations in four source
variables, with an invertible determinant. This criterion includes
flatness; reducedness of a special fibre alone would not suffice.
\end{proof}

\begin{theorem}\label{generic:degree}
For the declared marked localized source family there is a nonempty
base open $U\subset\Spec R$, containing $b_*$, such that
\[
 \overline X_U=X_U\longrightarrow U
 \quad\text{is finite \'etale of rank }512.
\]
In particular its characteristic-zero generic algebra over
$\mathbb Q(q_1,q_2,q_3)$ has dimension 512 and is \'etale. After
extension to an algebraic closure its generic fibre has exactly 512
reduced geometric points. The source is not required to be irreducible.
\end{theorem}
\begin{proof}
Let $O$ be the \'etale locus of $X$, an open subscheme of
$\overline X$. By Proposition~\ref{generic:controlfibre}, the entire
control fibre lies in $O$. The closed complement
$Z=\overline X\setminus O$ has closed image in the base because
$\overline X$ is proper. This image misses $b_*$. Remove it and call
the remaining nonempty open $U$. Then
$\overline X_U=O_U=X_U$, so this morphism is both proper and \'etale.
An \'etale morphism is locally quasi-finite; a proper locally
quasi-finite morphism is finite. A finite \'etale algebra is finite
locally free, with locally constant rank. The open $U$ is irreducible
because $R$ is a domain, and its rank is 512 at $b_*$. It therefore
has rank 512 everywhere on $U$, in particular at the characteristic-zero
generic point. This proves the theorem.
\end{proof}

The standard scheme-theoretic ingredients used here are restriction of
scheme closure to opens
(\href{https://stacks.math.columbia.edu/tag/01R5}{Stacks, Tag 01R5}),
the square Jacobian criterion
(\href{https://stacks.math.columbia.edu/tag/02GU}{Tag 02GU}),
proper and locally quasi-finite implies finite
(\href{https://stacks.math.columbia.edu/tag/02LS}{Tag 02LS}), and finite
\'etale local freeness
(\href{https://stacks.math.columbia.edu/tag/02K9}{Tag 02K9}). The finite
degree conclusion concerns this precise localization and marking; it
asserts no unmarked quotient, full moduli classification or physical
multiplicity.

\section{A uniform length bound at every finite fibre}
\label{length:section}

The generic degree computed in the preceding section is not by itself an
upper bound at a special metric. The following argument supplies precisely
the missing specialization statement. It applies to the original square
presentation, not to a determinantal plane quotient.

\begin{lemma}[Finite fibres of a square family]\label{length:bound}
Let $A$ be a Noetherian integral $\mathbb Q$-algebra, $K=\operatorname{Frac}(A)$,
and
\[
 B=A[x_1,\ldots,x_n,1/\Delta]/(f_1,\ldots,f_n),
 \qquad f_i,\Delta\in A[x_1,\ldots,x_n].
\]
Suppose $\dim_K(B\otimes_AK)=d<\infty$. If the entire fibre over
$s\in\Spec A$ is finite, then
\[
 \dim_{\kappa(s)}(B\otimes_A\kappa(s))\le d.
\]
The statement includes all multiplicities and does not assume properness,
reducedness of the special fibre, or that $B$ is finite as an $A$-module.
\end{lemma}
\begin{proof}
The empty fibre satisfies the assertion. Otherwise write its finitely many
support points as primes $\mathfrak p_1,\ldots,\mathfrak p_k$ of $B$.
They lie in the quasi-finite locus $U$ of $\Spec B\to\Spec A$. This locus
is open by the quasi-finite-locus theorem \cite{stacks-qf}. Write the
complement as $V(J)$. For every $i$, $J\not\subset\mathfrak p_i$.
Finite prime avoidance gives $h\in J$ outside all the $\mathfrak p_i$.
Consequently $D(h)\subset U$ and $D(h)$ contains the whole special fibre.
The image of $h$ is outside every maximal ideal of its Artinian coordinate
algebra; it is therefore a unit. Localizing at $h$ preserves that entire
fibre as a scheme, including its nilpotents.

Choose a polynomial representative $H$ with $h=H/\Delta^j$ in $B$.
Then, canonically,
\[
 B_h=A[x_1,\ldots,x_n,1/(\Delta H)]/(f_1,\ldots,f_n).
\]
Its morphism to $\Spec A$ is locally quasi-finite and has a localized
square polynomial presentation. The square-presentation criterion for
quasi-finite syntomic morphisms \cite{stacks-square}, condition (4)
implying (2), says that this morphism is flat. This use of flatness does
not assume the base regular. Introducing one inverse variable replaces
the presentation by $n+1$ variables and $n+1$ equations and gives the
same criterion.

The remaining argument bounds fibre dimension for a flat module that
need not be finitely generated. Let $M=B_h$ as an $A$-module. Its generic fibre is a localization of the
finite-dimensional algebra $B\otimes_AK$, so it has dimension at most $d$.
No finiteness of $M$ as an $A$-module has been inferred. Put $N=d+1$.
Since $N!$ is invertible in $A$, the antisymmetrizer
\[
 e=\frac1{N!}\sum_{\pi\in S_N}\operatorname{sgn}(\pi)\pi
       \quad\text{on }M^{\otimes_A N}
\]
is an idempotent whose image is canonically $\bigwedge_A^N M$.
Indeed, it annihilates repeated arguments and descends through the
exterior quotient, whose composition with the descended map is the
identity. Thus this exterior power is a direct summand of a flat module
and is flat. Over the domain $A$ it is consequently torsion-free.
Its localization to $K$ is zero, since $\dim_K(M\otimes_AK)\le d$.
A torsion-free module over a domain injects into its generic localization,
so $\bigwedge_A^N M=0$.

Exterior powers commute with base change. Therefore
\[
 \bigwedge_{\kappa(s)}^N(M\otimes_A\kappa(s))=0,
\]
which forces the vector-space dimension to be at most $N-1=d$.
The localization at $h$ preserved the full original special fibre,
so this is the required length bound.
\end{proof}

\begin{corollary}\label{length:application}
Every finite fibre of the declared localized connection family in
characteristic zero has scheme length at most $512$.
\end{corollary}
\begin{proof}
After fixing the four invertible scales, use the integral metric base
\[
 A=\mathbb Q[q_1,q_2,q_3,(q_1q_2q_3\prod_i(q_i-1))^{-1}].
\]
The scheme isomorphisms in Section~\ref{geometry:model} give four source
variables $(u_1,u_2,u_3,m)$ and the four rational equations
\[
 m^2+S-1=0,\quad AB+4z_2=0,\quad
 ABm+A(z_2+z_3)+B(z_1+z_2)=0,\quad H_0+mH_1=0.
\]
Here all quantities are the explicit functions defined in that section.
Multiply each equation by its denominator and combine those denominators
with the numerators of all inverted factors into a single $\Delta$.
On $D(\Delta)$ every multiplier is a unit. Hence this gives a square
presentation of exactly the same scheme, not only the same field-valued
points. Its generic algebra has dimension $512$ by
Theorem~\ref{generic:degree}. Lemma~\ref{length:bound}, with $n=4$,
proves the assertion. The same argument after base extension to
$\mathbb C$ gives the geometric-fibre statement.
\end{proof}

For clarity, a finite fibre has not yet been proved to have length
\emph{equal} to $512$. For example, the square family
$\mathbb Q[t,x]/(tx^2+x)$ has generic length two and special length one
at $t=0$. The exact equality below therefore needs both finiteness and
512 distinct lifted points. Neither a generic degree nor a list of
modular points is silently substituted for these two obligations.

\section{The exact fibre at a rational metric}\label{centre:section}
We now specialize to
\[
 q^0=(4/3,12/11,3/4)
\]
and keep the four scale parameters equal to one. Let $C$ be the rational
coordinate algebra of the original marked localized coefficient scheme at
this metric. Thus every original open condition remains inverted; $C$ is
not the algebra of the unfiltered projection locus. We prove finiteness
first, before applying the length bound of Lemma~\ref{length:bound}. This
order prevents a generic degree from being used as an unjustified bound
on an infinite exceptional fibre.

All large integer polynomials used below are supplied in full, with
ascending univariate coefficient arrays and expanded multivariate integer
coefficients, in \path{certificates/centre/}. The descriptions below define
the exact finite identity tests; no numerical root approximation or prior
review verdict is needed. In particular the two source minors, the
eliminant, its two factors, the reconstruction cofactors and all modular
coordinate residues are actual delivered polynomials.

\subsection{A two-generator plane module proves finiteness}
Use the scalar elimination and inverse reconstruction established in the
source reduction, and write $y=u_2$, $z=u_3$, $v=1/u_1$. Denote the three
primitive integer numerator polynomials after eliminating $m$ by
\[
 P_p,P_F,P_T\in\mathbb Z[u_1,y,z].
\]
Their degrees in $u_1$ are $2,6,10$, respectively. In terms of the source
quantities $\zeta_j=1-(u_j-1)^2/q_j$, the eliminated scalar is
\[
 m_0=\frac{A(\zeta_2+\zeta_3)+B(\zeta_1+\zeta_2)}{4\zeta_2},\quad
 A=\zeta_1u_2/u_1+\zeta_2u_1/u_2,\quad
 B=\zeta_2u_3/u_2+\zeta_3u_2/u_3.
\]
Here $\zeta_2$ is already a unit. The files \path{Hp.txt}, \path{F.txt}
and \path{Htriangle.txt} contain $P_p,P_F,P_T$, with the source meanings
specified by that reduction.

Define
\begin{align*}
 D&=27y^2z^2-32y^2z+4y^2-22yz^2-z^2,\\
 J&=11y^2-22y-1,\qquad K=20y^2-22y-1,\qquad E=3D+8zJ.
\end{align*}
For clarity, the exact source-to-file identities at this metric are
\begin{align}
 H_p&=\frac{P_p}{144u_1u_2^2u_3},\label{centre:source1}\\
 m_0^2+S-1&=\frac{P_F}{2304u_1^2u_2^2u_3^2J^2},\label{centre:source2}\\
 \mathcal H_0+m_0\mathcal H_1
 &=\frac{P_T}{20736u_1u_2u_3J
 (4u_3-1)^4(11u_2+1)^4(3u_1+1)^4}.\label{centre:source3}
\end{align}
These rational identities are checked by complete coefficient comparison
in Section~\ref{geometry:centre-source}. Their denominators are products
of declared units: $J=-11\alpha_2$, and the three linear factors are
$4\beta_3,11\beta_2,3\beta_1$. In these identities $m_0$ is the
displayed $C_p/(4\zeta_2)$, where
$C_p=A(\zeta_2+\zeta_3)+B(\zeta_1+\zeta_2)$. The alternative
$-C_p/(AB)$ agrees with it on $H_p=0$; it is not substituted off that
equation when deriving these numerator files.
The identities $D=-12yzB$ and $J=-12\zeta_2$ show that every original
point is on the open with coordinate ring
\[
 R=\mathbb Q[y,z,(yzDJ)^{-1}].
\]
Reversal in the invertible variable gives
\[
 \phi=v^2P_p(1/v,y,z)=a v^2+bv+c,\qquad
 a=-3y^2D,\quad b=-6y^2E,\quad c=KD.
\]
The leading coefficient $a$ is a unit of $R$. Consequently
\[
 P=\phi/a=v^2+\lambda v+\mu,\qquad
 \lambda=2E/D,\quad \mu=-K/(3y^2)
\]
defines a free rank-two $R$-algebra with basis $(1,v)$.

Put $U=v^6P_F(1/v,y,z)$ and $T_0=v^{10}P_T(1/v,y,z)$. Full integer
coefficient multiplication verifies
\begin{align}
 a^5U&=Q_F\phi+216y^6J^2(rv+s),\label{centre:pdiv1}\\
 a^9T_0&=Q_T\phi+1944y^{10}J(tv+w).\label{centre:pdiv2}
\end{align}
The quotients $Q_F,Q_T$ and the entire remainders are delivered, so these
are coefficient identities, not sampled equalities. The term counts of
$(r,s,t,w)$ are $(157,173,797,816)$; their respective bidegrees in $(y,z)$
are $(18,14),(18,14),(35,27),(35,27)$. All coefficients removed in
\eqref{centre:pdiv1}--\eqref{centre:pdiv2} are units of $R$. Hence the
ideals generated by $(\phi,U,T_0)$ and $(P,rv+s,tv+w)$ are equal in
$R[v]$. This equality holds over arbitrary $R$-algebras, including
nonreduced ones.

The candidate algebra
\[
 \mathcal B=R[v]/(P,rv+s,tv+w)
\]
has the exact module presentation
\begin{equation}\label{centre:module}
 R^4\xrightarrow{\mathcal K}R^2\longrightarrow\mathcal B\longrightarrow0,
 \qquad
 \mathcal K=\begin{pmatrix}
 s&-\mu r&w&-\mu t\\r&s-\lambda r&t&w-\lambda t
 \end{pmatrix}.
\end{equation}
Indeed, the four columns are the coordinates of
$(rv+s),v(rv+s),(tv+w),v(tv+w)$ in the basis $(1,v)$, and these generate
the ideal as an $R$-module. No division by $r$ or $t$ is used here.

Multiply columns two and four by $a$, a unit, and write
$\widehat{\mathcal K}=\mathcal K\operatorname{diag}(1,a,1,a)$.
Define $f,g\in\mathbb Z[y,z]$ by the exact normalization identities
\[
 \det\widehat{\mathcal K}_{12}=12y^2D^7f,\qquad
 \det\widehat{\mathcal K}_{23}=6y^2D^9g.
\]
The full polynomials are \path{minor_12.txt} and \path{minor_23.txt}.
The other four normalized minors are also supplied and checked; no factor
unsupported on an existing unit is discarded. The relevant degrees are
\[
 \deg_{y,z}f=(24,16),\quad \deg_{y,z}g=(37,25),\qquad
 \#\operatorname{terms}(f)=281,\quad\#\operatorname{terms}(g)=833.
\]

\begin{lemma}\label{centre:coprime}
The polynomials $f,g$ are coprime in $\mathbb C[y,z]$.
\end{lemma}
\begin{proof}
Exact extraction gives $[y^0]f=-z^{16}$ and
$([y^{24}]f)(0)=-40960000$. Thus $f$, considered in $\mathbb C[z][y]$,
has no nonconstant content factor. Write $F_1(y)=f(y,1)$ and
$G_1(y)=g(y,1)$. Their degrees are 24 and 37; their leading coefficients
modulo 101 are 4 and 83, respectively. An explicit extended-Euclidean
certificate, recomputed from these coefficient arrays, gives
\[
 a_*\overline F_1+b_*\overline G_1=1
 \quad\hbox{in }\F_{101}[y],\qquad
 \deg a_*=36,\quad\deg b_*=23.
\]
As a separate test the complete ascending-coefficient $61\times61$
Sylvester matrix has determinant $98$ modulo 101.

If a primitive nonconstant integer polynomial divided $F_1,G_1$, its
reduction would still have positive degree: reduction preserves the
entire degree of $F_1$, so it preserves the degree of each factor of
$F_1$. This contradicts the displayed identity. Thus $F_1,G_1$ are
coprime over $\mathbb Q$, and a rational B\'ezout identity proves the
same over $\mathbb C$. A common factor of $f,g$ of positive $y$-degree
would stay nonconstant at $z=1$, because that specialization preserves
the full $y$-degree of $f$: the leading $y$-coefficients of factors
multiply, and their product is nonzero at $z=1$. It would divide $F_1,G_1$. A common factor
of $y$-degree zero would divide the content of $f$. Both possibilities
have been excluded.
\end{proof}

\begin{theorem}\label{centre:finite}
The entire original localized fibre at $q^0$ is finite, including every
possible nonreduced contribution.
\end{theorem}
\begin{proof}
Two coprime polynomials in two variables have a zero-dimensional common
zero scheme: a positive-dimensional component would give a height-one
prime, which is principal in the polynomial ring, and hence a common
factor. Therefore $\mathbb Q[y,z]/(f,g)$, and its localization
$R/(f,g)$, are finite-dimensional vector spaces.

Every maximal minor of a two-row presentation matrix annihilates its
cokernel. For a selected pair of columns, their $2\times2$ adjugate
expresses the determinant times either basis vector in the image.
The two displayed minors differ from $f,g$ by units, so
$(f,g)\mathcal B=0$. Presentation~\eqref{centre:module} gives a
surjection $(R/(f,g))^2\twoheadrightarrow\mathcal B$. Thus
$\mathcal B$ is finite dimensional. The original coordinate algebra
$C$ is obtained by the inverse change of variables, the unique source
reconstruction and all remaining original localizations. Localization
of a finite-dimensional algebra selects Artinian factors and remains
finite dimensional. These are equalities of localized schemes; in
particular the argument preserves nilpotents. No rank-zero stratum or
additional localization boundary was discarded.
\end{proof}

Applying Lemma~\ref{length:bound} to the original square family, with
generic degree 512 and the finite fibre just proved, now gives
\begin{equation}\label{centre:upper}
 \dim_{\mathbb Q}C=\dim_{\mathbb C}(C\otimes\mathbb C)\le512.
\end{equation}
The equality follows from extension of scalars for a finite vector space.
This is the first point in the argument at which the generic degree is
used to bound this particular fibre.

\subsection{A complete smooth residue witness and its lifts}
The file \path{candidate_mod101.json} supplies eleven monic polynomials
$h_i(Z)$ and residues $U_{1i},U_{2i},U_{3i}$ in
\[
 T_i=\F_{101}[Z]/(h_i),\qquad U_{3i}=Z\pmod {h_i}.
\]
The certificate is fully checked in these quotient rings by polynomial
convolution, Euclidean division and explicit inverses. The degrees are
\begin{equation}\label{centre:degrees}
 (1,1,2,2,5,6,56,61,62,158,158),\qquad\sum_i\deg h_i=512.
\end{equation}
The tests are
\begin{align}
 \gcd(h_i,h_i')&=1,&\gcd(h_i,h_j)&=1\quad(i\ne j),\label{centre:squarefree}\\
 P_p(U_i)=P_F(U_i)=P_T(U_i)&=0,&
 \det\frac{\partial(P_p,P_F,P_T)}{\partial(u_1,u_2,u_3)}(U_i)&\in T_i^\times.
 \label{centre:smooth}
\end{align}
The Jacobian is that of the original three numerator equations, rather
than of the two convenient projection minors. The first two components
are $h_1=Z+48$, $U_1=(29,82,53)$ and $h_2=Z+67$,
$U_2=(47,39,34)$; the remaining, possibly reducible, component arrays
are supplied in full. Irreducibility is unnecessary:
\eqref{centre:squarefree} gives 512 distinct geometric values of $u_3$.

Every reconstruction denominator is a unit in each $T_i$. In particular,
with
\[
 \alpha_j=q_j-(u_j-1)^2,\quad\beta_j=q_j+u_j-1,\quad
 d_j=q_ju_j/\beta_j,\quad X_j=\alpha_j/\beta_j,\quad Y_j=\beta_j/q_j,
\]
the full list of original factors checked is
\begin{equation}\label{centre:23}
 \begin{gathered}
 d_j,\ q_j-d_j,\ q_j-d_j^2\quad(j=1,2,3),\\
 S,A,B,E_1,p_0,2-p_0,t_0,2-t_0,p,2-p,r_1,2-r_1,t_1,2-t_1,
 \end{gathered}
\end{equation}
where $S=\sum_j\zeta_j$, $E_1=\sum_jX_j^2(Y_{j+1}^2-Y_{j+2}^2)$
with cyclic indices, $p_0=2\zeta_2u_1/(u_2A)$ and
$t_0=2\zeta_3u_2/(u_3B)$. A unit test means a polynomial gcd equal to
one, equivalently an inverse identity modulo the entire $h_i$.
The triangle is reconstructed from its three right-kernel equations,
two left-kernel equations and its trace equation; the determinant is
$Y_3E_1$. The remaining left-kernel equation and all ten original
coefficient residuals are checked after reconstruction. The linear
remainder coefficient $r(U_{2i},U_{3i})$ is separately a unit.

\begin{lemma}\label{centre:hensel}
The original characteristic-zero fibre has 512 distinct points after
extension to a single finite extension of $\mathbb Q_{101}$, with 512
distinct values of $u_3$.
\end{lemma}
\begin{proof}
Choose a finite field $k/\F_{101}$ splitting all $h_i$, and let
$\mathcal O$ be the complete ring of integers of the corresponding
finite unramified extension $L/\mathbb Q_{101}$. Each residue tuple is
a common zero of the three integer source polynomials with invertible
Jacobian. Suppose $a_k\in\mathcal O^3$ solves these equations modulo
$101^k$. Choose a correction $t\pmod {101}$ solving
\[
 J(a_k)t\equiv-101^{-k}(P_p,P_F,P_T)(a_k)\pmod {101},
 \qquad a_{k+1}=a_k+101^kt.
\]
The invertible residue Jacobian makes this possible, and the polynomial
Taylor formula proves the equations modulo $101^{k+1}$. Completeness
gives an exact source zero in $\mathcal O^3$.

All required rational constants are defined at 101. Every denominator
and factor in \eqref{centre:23} remains a unit, since its residue was
checked nonzero. Source reconstruction therefore gives an $L$-point of
the original localized fibre. Distinct residue tuples have distinct
lifts. More strongly, their $u_3$ residues are all distinct by
\eqref{centre:squarefree}, so their lifted $u_3$ values are distinct.
Only this witnessed subset is lifted; no assertion that the entire
special fibre has been described is needed.
\end{proof}

\begin{theorem}\label{centre:complex}
The original marked localized coefficient fibre at $q^0$ has exactly
512 distinct reduced nonsingular complex points and scheme length 512:
\[
 C\otimes_{\mathbb Q}\mathbb C\cong\mathbb C^{512}.
\]
\end{theorem}
\begin{proof}
The 512 $L$-points of Lemma~\ref{centre:hensel} give 512 distinct
maximal ideals with residue field $L$ in $C\otimes L$. The Chinese
remainder theorem gives a surjection
\[
 C\otimes L\twoheadrightarrow L^{512}.
\]
The source has dimension at most 512 by \eqref{centre:upper}, so the
surjection is an isomorphism and $\dim_{\mathbb Q}C=512$.
An additional point or a nonzero nilpotent contribution would add positive
length, so both are excluded by this equality. Faithful
field extension implies that $C$ is reduced. A finite reduced rational
algebra is a product of finite separable number fields; after extension
to $\mathbb C$ these split into exactly 512 copies of $\mathbb C$.
The original three-source Jacobian is nonzero at every lift by its
nonzero residue, so its determinant is a unit in $C\otimes L$ and
in $C$. The source equivalence preserves this nonsingularity.
This argument does not identify $\overline{\mathbb Q}_{101}$ with
$\mathbb C$, assume properness, or deduce characteristic-zero equality
from a modular count alone.
\end{proof}

\begin{corollary}\label{centre:primitive}
The coordinate $u_3$ generates $C$ as a rational algebra, and the
coefficient $r(u_2,u_3)$ is a unit in $C$.
\end{corollary}
\begin{proof}
In $C\otimes L=L^{512}$, the 512 entries of $u_3$ are distinct.
Vandermonde interpolation gives $L[u_3]=L^{512}$. The inclusion
$\mathbb Q[u_3]\subset C$ becomes an equality after faithful extension
to $L$, hence was an equality already. At every lift, $r$ has nonzero
residue. Thus multiplication by $r$ is invertible on $C\otimes L$;
its determinant is then nonzero over $\mathbb Q$, so $r$ is a unit in
$C$. The relation $rv+s=0$ gives $v=-s/r$ everywhere, including every
point of the original fibre.
\end{proof}

\subsection{Identification of the explicit rational univariate algebra}
The two given source minors define the integer resultant
\[
 R_*(z)=\operatorname{Res}_y(f(y,z),g(y,z)),\qquad
 \deg R_*\le24\cdot25+37\cdot16=1192.
\]
The exact supplied coefficient arrays satisfy
\begin{equation}\label{centre:factorization}
 R_*=c_*HE_*,\quad c_*\in\mathbb Z\setminus\{0\},\quad
 \deg R_*=1192,\quad\deg H=512,\quad\deg E_*=680,
\end{equation}
with $H,E_*$ primitive. This complete product identity is checked by
integer convolution. The source resultant identity is checked by a
separate coefficient parser and exact integer univariate resultants at
1193 distinct negative integers. At every test the $y$-degrees remain
24 and 37. Both the actual resultant and the proposed array have degree
at most 1192, so these tests prove equality of polynomials over
$\mathbb Z$. These are a complete identity proof, rather than random
spot checks. FLINT supplies the exact integer resultant backend;
the product and modular factor tests use only Python integers.

The complete complementary factor, with every multiplicity retained,
satisfies
\begin{equation}\label{centre:Hmod}
 \operatorname{lc}(H)\equiv10\pmod {101},\quad
 \overline H/\overline{\operatorname{lc}(H)}=\prod_i h_i,
 \qquad \gcd(\overline E_*,h_i)=1\quad\hbox{for every }i.
\end{equation}
\begin{theorem}\label{centre:coordinate}
The supplied polynomial $H$, after monic normalization, is the minimal
polynomial of $u_3$ in the original rational coordinate algebra:
\[
 C=\mathbb Q[u_3]\cong\mathbb Q[Z]/(H(Z)).
\]
In particular this algebra is identified in characteristic zero, not
merely by agreement of its reduction modulo 101.
\end{theorem}
\begin{proof}
At each of the 512 integral lifts, the source minors $f,g$ vanish.
The Sylvester adjugate identity therefore gives $R_*(u_3)=0$, even if
a leading coefficient of $f$ or $g$ drops at that point. By
\eqref{centre:Hmod}, $E_*(u_3)$ has nonzero residue and is a unit in
$\mathcal O$. The exact identity \eqref{centre:factorization} and
$c_*\ne0$ in the characteristic-zero field $L$ imply $H(u_3)=0$.
The scalar $c_*$ need not be a 101-adic unit; it is cancelled in $L$.
Thus $H$ annihilates $u_3$ in $C\otimes L$ and, by faithful base
change, in $C$. Corollary~\ref{centre:primitive} shows that its minimal
polynomial has degree 512. The displayed $H$ has the same degree, so
its monic normalization is exactly that polynomial.
\end{proof}

\subsection{A complete exact real-root certificate}
Write $H(x)=\sum_{i=0}^{512}a_ix^i$. The coefficient arrays satisfy
\[
 \sum_{i<512}|a_i|B^i<|a_{512}|B^{512},\qquad B=2^{19}.
\]
For $|x|\ge B$, each ratio $|a_ix^i|/|a_{512}x^{512}|$ is no larger
than at $B$. Thus all complex roots have modulus strictly less than
$B$. In particular every real root is in $(-B,B)$.

Parameterize the negative and positive halves by $x=-Bt$ and $x=Bt$,
$0<t<1$. For a node polynomial $Q(t)$ of degree $n$, define
\[
 \mathcal DQ(s)=(1+s)^nQ(1/(1+s)).
\]
For $s>0$, this change of variables bijects to $0<t<1$; its prefactor
never vanishes there. Descartes' rule of signs bounds the positive-root
multiplicity count by the variation $V$ of the nonzero coefficients
of $\mathcal DQ$, and their difference is even. In particular $V=0$
proves emptiness and $V=1$ proves exactly one root. Coefficient reversal
followed by a Taylor shift by one computes $\mathcal DQ$ exactly.

For $V>1$ bisect, using the integer coefficient transformations
\[
 Q_0(t)=2^nQ(t/2),\qquad Q_1(t)=Q_0(t+1).
\]
Positive integer content can be removed without changing roots or
signs. A root at a subdivision midpoint is tested exactly and would
be counted once, separately from the two open children. The delivered
finite trees, in \path{descartes_partition.json}, have the following
independently checked data:
\begin{center}
\begin{tabular}{lrrrr}
\toprule
Half-axis & One-root leaves & All leaves & Nodes & Maximum depth\\
\midrule
$u_3<0$ &46&85&169&38\\
$u_3>0$ &158&233&465&35\\
\midrule
Total &204&318&634&\\
\bottomrule
\end{tabular}
\end{center}
Every leaf is checked to have variation zero or one. The leaf paths
are distinct and prefix-free, and every internal prefix has both
children. Therefore the finite trees cover the initial intervals,
apart from the separately tested splitting points. None of the
splitting points is a root, and $H(0)\ne0$. A second integer-only
checker independently verifies that all 319 distinct partition
endpoints are nonroots and that each of the 204 one-root leaves has
opposite endpoint signs. This also supplies a direct intermediate-value
lower bound. No assumed termination of an infinite subdivision is
used: the delivered finite tree is the complete certificate.

\begin{theorem}\label{centre:real}
The original localized real coefficient fibre has exactly 204 points:
\[
 N_{\mathbb R}=204,
 \qquad N_{\mathbb R,u_3<0}=46,\quad N_{\mathbb R,u_3>0}=158.
\]
The remaining 308 points of the complexification form 154 conjugate
pairs. The 204 real coefficient points give the real, star-preserving
connections in the declared sector through the source equivalence.
\end{theorem}
\begin{proof}
The strict root bound, full binary coverage, endpoint tests and Descartes
variations prove that $H$ has exactly $46+158=204$ real roots. By
Theorem~\ref{centre:coordinate}, real points are precisely the rational
algebra homomorphisms $C\to\mathbb R$, which are exactly evaluations
at those roots. Every original coordinate and every localized inverse
belongs to $C=\mathbb Q[u_3]$, so it evaluates to a real value; no
additional reality filter is missing. The geometric identification
with star-preserving connections is the real source equivalence, not
a claim that arbitrary complex coefficients satisfy an anti-linear
star condition. Theorem~\ref{centre:complex} and complex conjugation
give $(512-204)/2=154$ nonreal pairs.
\end{proof}

\subsection{An explicit global reconstruction, with safe cancellation}
For completeness the other coordinates can be reconstructed without
an unspecified interpolation choice. Form a $59\times60$ coefficient
matrix $V(Z)$, with columns indexed by $1,y,\ldots,y^{59}$ and rows
\[
 f,yf,\ldots,y^{35}f,\qquad g,yg,\ldots,y^{22}g.
\]
Let $V_j$ denote its degree-$j$ column and set
\[
 A_*(Z)=\det[V_1|V_2|\cdots|V_{59}],\qquad
 B_*(Z)=\det[V_0|V_2|\cdots|V_{59}].
\]
If $c_i=(-1)^i\det(V_{\ge2}\text{ with row }i\text{ removed})$,
expanding along the first column gives the constant and linear
coefficients $B_*,A_*$. Every higher coefficient in that row
combination is zero by a repeated-column determinant. Consequently
\begin{equation}\label{centre:cofactorrel}
 A_*(Z)y+B_*(Z)\in(f,g)\mathbb Z[y,Z].
\end{equation}
The degree bound is $36\cdot16+23\cdot25=1151$. All coefficients of
both determinants are supplied; equality is checked at 1152 distinct
negative integers using the entire fixed-size matrices. Degree drops
of specialized input polynomials do not affect these fixed
determinant identities.

The exact integer arrays satisfy
\[
 A_*=c_0GD_2,\qquad B_*=-c_0GN,
 \qquad c_0\in\mathbb Z\setminus\{0\},
 \quad \deg G=388,\quad\deg N=\deg D_2=763.
\]
Integer convolution checks both identities. Full component tests give
$\gcd(\overline G,h_i)=\gcd(\overline {D_2},h_i)=1$ for every $i$.
At all 512 integral lifts, $G,D_2$ therefore have nonzero residue.
The lifts exhaust $C\otimes L$, so both are units there and descend
to units of $C$. In $C$, the source minors vanish; cancelling these
proved units in \eqref{centre:cofactorrel} gives
\[
 u_3=Z,\qquad u_2=N(Z)/D_2(Z),\qquad
 u_1=-r(u_2,Z)/s(u_2,Z).
\]
Here $s=-rv=-r/u_1$ is a unit because $r$ and $u_1$ are units.
Formal cancellation in $\mathbb Q(Z)$ alone would not have justified
these formulas at the fibre points.

The remaining coordinates are given by $m=m_0$, the displayed path
formulas, and the unique triangle linear solve. Explicitly, in the
unknown order $(a,b,c,p,r_1,t_1)$ the matrix and right-hand side are
\[
 M=\begin{pmatrix}
 Y_1&0&0&Y_2&Y_3&0\\
 0&Y_2&0&-Y_1&0&Y_3\\
 0&0&Y_3&0&-Y_1&-Y_2\\
 X_1&0&0&-X_2&-X_3&0\\
 0&X_2&0&X_1&0&-X_3\\
 1&1&1&0&0&0
 \end{pmatrix},\quad
 w=\begin{pmatrix}
 -mY_1\\-mY_2-2Y_1\\-mY_3-2Y_1-2Y_2\\
 -mX_1-2X_2-2X_3\\-mX_2-2X_3\\-m
 \end{pmatrix},\qquad \det M=Y_3E_1.
\]
Cramer's rule gives a literal rational circuit. Its complete 53-node,
30-division serialization is \path{rational_circuit.json}; its inputs
$N,D_2$ are \path{u2_reconstruction.json}. Every division denominator
is checked to have an explicit inverse in every $T_i$. Inductively,
earlier unit divisions preserve integrality at every lift, and the
next nonzero residue proves its denominator is a unit there. Since
the lifts exhaust the algebra, each denominator is a unit in $C$.
The exact cofactor relation, the source linear remainder and the unique
source reconstruction prove characteristic-zero equality with the
original coordinates. Modular residual tests corroborate this
implementation; they are not used as a substitute for those exact
identifications.

At each real root of $H$ the circuit therefore produces one and only
one original real coefficient point, with all 23 factors nonzero.
Conversely every original point is obtained in this way. The
construction counts marked coefficient points in this localized
sector; it does not divide by relabellings or classify a larger moduli
space.

\paragraph{A worked real connection.}
At $q^0$ with all four scales one, let $\rho$ be the unique root of the
primitive polynomial $H$ in
\[
 -\frac{400542}{10^6}<\rho<-\frac{400541}{10^6}.
\]
The exact $H$-values at the endpoints have negative and positive signs,
respectively, and this
interval lies inside the certified one-root Descartes leaf
$(-1/2,-1/4)$.  Set $u_3=\rho$ and use the existing rational
reconstruction circuit.  The spatial blocks $M_0,M_1$ from the original connection ansatz and
the scalar have the following certified rounded values:
\[
\begin{aligned}
 (u_1,u_2,u_3)&\simeq (-0.4187284043,-0.4102925777,-0.4005416805),\\
 m&\simeq -1.9869902124,\\[-0.2em]
\end{aligned}
\]
\[
 M_0\simeq\begin{pmatrix}0.757841 & 1.254421 & 0 \\ 0.745579 & -0.088125 & 1.346200 \\ 0 & 0.653800 & 1.317274\end{pmatrix},\qquad
 M_1\simeq\begin{pmatrix}0.974629 & 0.528544 & -0.103648 \\ 1.471456 & -0.376875 & 0.689222 \\ 2.103648 & 1.310778 & 1.389236\end{pmatrix}.
\]
For example, the temporal mixed coefficients are
\[
\begin{aligned}
 (v_{01},v_{02},v_{03})(0)&\simeq (7.956630,2.811717,1.862320),\\
 (v_{01},v_{02},v_{03})(1)&\simeq (1.217001,2.006328,4.032929).
\end{aligned}
\]
An exact nested dyadic bracket and outward-rounded interval evaluation
of the literal circuit certify all displayed rounding.  Every one of
the original twenty-three units has absolute value greater than
$0.1036$ on that bracket.  The exact algebraic connection is specified by $\rho$ and the rational
circuit. These rounded displays illustrate its active spatial and temporal
couplings. The supplementary folder \path{examples/worked_connection/}
contains the exact signs, outward enclosures and a separate replay script.

\section{Closure of the main theorem}
\begin{proof}[Proof of Theorem~\ref{main:count}]
The scheme-theoretic reconstruction of Section~\ref{geometry:model}
identifies the original declared connection sector with the four-variable
family, preserving all inverted factors and all multiplicities.
Theorem~\ref{generic:degree} proves the first assertion by a proper
closure, complete boundary exclusion at a control point and the actual
relative Jacobian criterion.

At $q^0$, Theorem~\ref{centre:finite} proves that the entire localized
fibre is finite. Corollary~\ref{length:application} therefore bounds its
length by $512$. The modular smooth-point construction and its
characteristic-zero lifts in Theorem~\ref{centre:complex} supply $512$
different points of that same fibre. A point of an Artinian scheme
over an algebraically closed field contributes length at least one.
Thus the length cap is attained, every point contributes length one,
and no extra support or nilpotent contribution remains. This proves
the second assertion.

The exact coordinate-algebra identification and Descartes certificate
in Theorem~\ref{centre:real} prove that precisely $204$ of these points
are real. The rational inverse reconstruction is defined at each point
and is over $\mathbb Q$, so real primitive coordinates give real
coefficient points and conversely. The original real QLC isomorphism
then supplies the claimed star-preserving connections. This proves
the third assertion.
\end{proof}

% END SHARED COUNTING EXPOSITION
\section{The computational certificates and reproducibility}
\label{verification:section}

\subsection{What the computer-assisted part proves}
The proof uses explicit finite certificates, grouped by mathematical role.
The directory \path{certificates/geometry/} supplies the original global
equations and all cleared localization factors. Its checking programs
derive the tensor-component tables, the ring identities and the path and
triangle completion identities, and compare the full coefficient arrays
with the intrinsic formulas in Section~\ref{geometry:model}.

The additional direct checker \path{audit/qlc_original_replay.py}
imports no other checking module. Starting with symbolic temporal
coordinates, all four metric scales and the supported spatial matrices,
it reconstructs all $56$ coefficient values and expands the original
tensor connection and conjugate tensor flip. It checks the explicit
ideal representations \eqref{integrity:direct-residuals} for all $128$
metric and $128$ star components, torsion, degree preservation, and
both compositions needed for $\sigma^{-1}=D\sigma D$.
It also checks the right Leibniz rule on every vertex idempotent at
every vertex: $16\cdot16\cdot4\cdot4\cdot4=16384$ scalar equalities.
Two exact negative controls change a reconstructed square coefficient
or break an active edge sum; their residual discrepancies are nonzero
rational functions. No random or floating-point point test is used.
This is a direct check of inverse reconstruction, conditional on the
block involution ideal. The separate block certificate proves that
the four reduced equations imply that ideal, and the universal-ring
argument proves the forward elimination and the identity of the two
maps. Neither direction is inferred from the other by a point count.

The directory \path{certificates/generic/} supplies finite-field
coefficients and unit identities for every finite boundary chart and
pole chart, the restricted elimination data, the $512$-degree control
algebra, and the relative Jacobian tests. Their role is to check the
finite algebraic calculations in the generic-degree proof; the
properness and specialization arguments are proved in the text.

The directory \path{certificates/centre/} contains all numerator
polynomials used in the quadratic-module reduction, all six minors,
the eleven smooth residue components, the full integer eliminant and
its factorization, the primitive coordinate polynomial, the complete
Descartes partition, and the rational reconstruction circuit. These are
the actual coefficient arrays, not only their checksums.
Every division that occurs in a coordinate reconstruction is justified
by a unit condition or an explicit inverse certificate. The proof does
not equate a determinantal plane quotient with the original connection
scheme, and does not discard a rank-zero locus without an argument.

For a long polynomial identity over $\mathbb Z$, exact evaluation and
interpolation are a proof when the declared degree bound is justified
and the number of distinct exact test points exceeds that bound.
The eliminant and cofactor checks use grids disjoint from the production
grids. A handful of additional determinant samples is only
corroboration. Finite-field unit statements are verified by explicit
products equal to one, or B\'ezout identities, in the specified quotient
algebra. No floating-point residual is used to prove an exact equality.

\subsection{How to read or replay the proof}
The package root contains \path{manuscript.tex}, \path{manuscript.pdf},
the complete standalone \path{supplement.tex} and \path{supplement.pdf},
\path{audit/}, \path{certificates/}, \path{requirements.txt} and the portable
worked example in \path{examples/worked_connection/}. The common counting
exposition and geometric consequences are identical in both TeX sources,
including the text between formal statements and proofs. The files in
\path{sections/} are retained literal source-binding inputs, not verbatim
extracts of the current printed sections; the complete documents are the
canonical exposition. The checker
\path{check_consistency.py} compares these complete regions, checks all
references and cited package paths, and binds both sources to this release.
No separate historical counting-proof document is required.

After installing \path{requirements.txt}, run
\begin{verbatim}
python -B reproduce.py --output ../qlc-results
\end{verbatim}
The output directory must be new and outside the package. The entry point
first checks source and package consistency, then runs \path{verify.py}
in both normal and explicit Python \texttt{-O} modes, and finally replays
the worked example with outward-rounded interval arithmetic. The counting
runner stages
only $175$ catalogued proof inputs, including $142$ coefficient files,
in a fresh temporary workspace. Prior outputs and review reports are
never proof inputs. Each child must exit successfully and create a new
receipt with its exact full-scope result status. The complete determinant
checks cannot be replaced by supplemental samples.

The supplement reports complete normal and optimized executions, their
receipt hashes, and the bytewise comparison of all mathematical result
pairs. The stated environment is Python~3.12.14, SymPy~1.14.0,
python-flint~0.9.0 and native FLINT~3.6.0.

An exact cross-block join binds all nine graph equations, all twenty-three
localization factors and both triangle invariants across the source and
degree calculations. Additional literal-formula checks bind the rational
circuit and the compact-top scalar normalization. With
$T_0=-P_{\rm top}|_{q_*}$, equation~\eqref{generic:compacttop} has scalar $-4$;
using the raw top polynomial instead has scalar $+4$. Both scalar
identities are checked in characteristic zero.

Python integers supply direct polynomial and quotient arithmetic.
SymPy and FLINT are explicit arithmetic dependencies for the specified
normalizations and large exact determinants. Hashes identify coefficient
bytes; executed identities and the written arguments establish the
mathematics. Scheme isomorphisms, properness, the finite-fibre length
lemma, Hensel lifting and faithful descent are proved in the text and
are not inferred from a software status label.

\subsection{Use of generative AI in the research workflow}
ChatGPT and OpenAI Codex assisted exploration of the algebraic reductions,
proof drafting, literature triage, development of checking programs and
internal adversarial reviews. The deterministic certificate checks use the
explicit coefficient arrays and algorithms described above. Their input
catalogue, exact outputs and executable replay are supplied in full, so
independent reproduction requires the stated arithmetic environment and
does not require a generative model or an earlier conversation.
The scheme identification, properness, length and lifting arguments are
written mathematical arguments requiring critical review. Internal agent
reviews contributed to checking the preparation; they are separate from
the exact arithmetic certificates and from external journal refereeing.

\subsection{Scope of the theorem and of the review}
The family is restricted by its chosen calculus, spatial homogeneity,
ordered path and triangle, trace constraints and the exact regular
open. These conditions are part of the definition. They have not been
proved to exhaust all connections on the cube. Nor is the number
divided by a relabelling group or interpreted as a count of physically
inequivalent geometries.

The result is the consolidated degree and fixed-metric count.
No explicit parameter-radius optimization, quantum functional integral,
regulator removal or physical continuum claim is needed for it.
The number $204$ is obtained only after the complexified coordinate
algebra and its real structure have both been identified.

\section{Conclusion}
The declared path--triangle QLC family admits an exact reversible
algebraic model. The corrected boundary and degree argument, finite-fibre
length bound and smooth-lift exhaustion establish a complete reduced
$512$-point fibre, not merely a modular lower bound. Its primitive
coordinate algebra and exact real-root certificate give $204$ actual
star-preserving connections at the stated rational metric.

% BEGIN SHARED GEOMETRIC CONSEQUENCES
\begin{corollary}[Distinct tensor Leibniz maps at one metric]
At $q^0$ with scales one, the $204$ counted real connections have $204$
distinct tensor Leibniz maps $\sigma$.
\end{corollary}
\begin{proof}
For this calculus the translation-degree argument in
Section~\ref{geometry:model} forces the inner correction map $\gamma$
to vanish. Hence
$\nabla\omega=\theta\otimes\omega-\sigma(\omega\otimes\theta)$.
Two equal maps $\sigma$ would give equal connections and equal coefficient
points, contrary to the distinctness already proved.
\end{proof}

This observation locates the finite multiplicity in the varying
bimodule data; it does not conflict with a comparison fixing the same
$\sigma$. Classification of excluded support or denominator strata,
special metrics and symmetry quotients are further questions. The
enumeration supplies no curvature formula, action principle or physical
selection of this support.

\begin{corollary}[Obstruction to ordinary coordinate products]
\label{geometry:noproduct}
None of the $204$ real connections counted at $q^0$ is an ordinary
tensor product of connections for a nontrivial partition of the four
coordinate directions. This remains true after permuting those direction
labels.
\end{corollary}
\begin{proof}
For an ordinary product over a partition, each frame element $e^i$
belongs to one factor and $\nabla e^i$ has both output directions in
that factor. Therefore $C^i{}_{jk}=0$ whenever $i,j,k$ do not all lie
in a single block. Take the block containing direction $0$. A nontrivial
partition has some spatial direction $a$ in another block, so it would
require $C^0{}_{0a}=v_{0a}=0$. On our declared open, however,
\eqref{geometry:temporal} gives
\[
 v_{0a}(0)=\frac{q_a-d_a^2}{q_a-d_a},\qquad
 v_{0a}(1)=\frac{q_a-d_a^2}{d_a(q_a-d_a)},
\]
and every displayed numerator and denominator is a unit. These values
cannot be zero at a real point. The three active pairs $\{0,a\}$
form a connected graph on all four directions; connectedness and the
same obstruction survive any permutation of the labels.
\end{proof}

The statement concerns ordinary products of the specified coordinate
calculi. It excludes neither a more general extension construction nor
an alternative change of geometric data.

% END SHARED GEOMETRIC CONSEQUENCES
\clearpage
\appendix
% BEGIN SHARED AW COMPARISON
\section{An explicit comparison with Aschieri--Weber compatibility}
\label{aw:appendix}
This appendix supplies the tensor calculation behind the same-metric
comparison with \cite[Definition~5.3(ii), Theorem~5.7]{orig-aw25}.
Its conventions and hypotheses are retained: it is not a comparison
obtained merely by fixing our unknown tensor Leibniz map $\sigma$.

Put $A_G=\mathbb C(G)$, $G=(\mathbb Z/2)^4$, and $k_i=h_i^{-1}$.
Edge symmetry gives $R_ik_i=k_i$. For the right-linear dual frame
$\xi_i$, with $\langle\xi_i,e^j\rangle=\delta_{ij}$, the corresponding
inverse metric is $\mathfrak g=\sum_i k_i\xi_i\otimes_{A_G}\xi_i$.
This calculus is bicovariant; the canonical Woronowicz map $\sigma_W$
flips its invariant frame tensors, and the exterior relations are
Grassmann relations. This canonical map is distinct from the variable
$\sigma$ of our bimodule connections.

\paragraph{The rational tensor product.}
Let $\phi_0$ be a right-linear coefficient map constant on invariant
frames and let $\phi=f\phi_0$, for any $f\in A_G$. Write $S_H$ for
the Hopf antipode and use Sweedler indices for the iterated coproduct
of $f$. The internal-Hom coaction and rational tensor formula
\cite[equations~(2.2), (2.6), (2.11)]{orig-aw25}, on invariant inputs, give
\[
 (\mathrm{id}\otimes_{\sigma_W}\phi)(e^i\otimes_{A_G}e^j)
 =\sum f_{(1)}e^i\otimes_{A_G}
             S_H(f_{(2)})f_{(3)}\phi_0(e^j).
\]
Moving the second coefficient through $e^i$ yields the coefficient
\[
 \sum f_{(1)}R_i\!\left(S_H(f_{(2)})f_{(3)}\right)=f,
 \qquad
 \sum f_{(1)}\otimes S_H(f_{(2)})f_{(3)}=f\otimes1.
\]
The second equality is the antipode identity and proves the first.
Thus the internal-Hom contribution and tensor balancing cancel the
apparent coefficient shift for every vertex function, without a
covariance or spatial-homogeneity restriction on $f$. For connections,
the rational tensor product in \cite[equation~(2.23), Section~4.4]{orig-aw25}
reduces on invariant frames to the same calculation: split off the
canonical connection, which vanishes on those frames and their tensors,
and apply the calculation to the remaining right-linear map.

Contracting with $\mathfrak g$ on
$(e^i\otimes e^j+e^j\otimes e^i)/2$ now gives the actual specialization
of \cite[equations~(5.9)--(5.10)]{orig-aw25}:
\[
 M_{ijb}=\Phi_{\mathfrak g}(L)_{ijb}
          =k_jL^i{}_{jb}+k_iL^j{}_{ib},\qquad
 L^i{}_{jb}=L^i{}_{bj},\quad M_{ijb}=M_{jib}.
\]
Its inverse is
\[
 L^i{}_{jb}
 =\frac{k_bM_{ijb}+k_jM_{ibj}-k_iM_{jbi}}{2k_jk_b}.
\]
Substitution verifies both compositions; the permuted cross terms
cancel and the remaining numerator is $2k_jk_bL^i{}_{jb}$.
All denominators are metric units. Consequently $\Phi_{\mathfrak g}$
is invertible here, including on arbitrary coefficient functions, and
the hypotheses of their Theorem~5.7 apply.

\paragraph{The compatibility predicate and the left/right passage.}
Write $\nabla_Re^i=\sum_{a,b}K^i_{ab}e^a\otimes_{A_G}e^b$.
The same tensor calculation specializes their metric condition to
\begin{equation}\label{aw:metric}
 k_jK^i_{jb}+k_iK^j_{ib}
       =\delta_{ij}(R_bk_i-k_i).
\end{equation}
Under $\nabla_R=D\circ\nabla\circ *$, the star and tensor flip of
Section~\ref{geometry:model} give
\[
 K^i_{ab}=-R_aR_b\bigl(\overline{C^i{}_{ba}}\bigr),
 \qquad K^i_{ii}=-\overline{U_i}.
\]
Setting $i=j=b$ in \eqref{aw:metric}, and using $R_ik_i=k_i$,
forces $U_i=0$ for all four directions. This conflicts with the
imposed trace $\sum_iU_i(t)=4$ in \eqref{geometry:support}.

The difference also persists without that trace restriction. Put
$Q_i(t)=h_i(1-t)/h_i(t)$ for $i>0$. The inverse formula above gives
\[
 K^i_{i0}(t)=\frac{k_i(1-t)-k_i(t)}{2k_i(t)}
            =\frac{Q_i(t)^{-1}-1}{2},\qquad K^0_{0i}=0.
\]
Pulling back shifts the first coefficient to the opposite time slice;
hence $v_{i0}=(1-Q_i)/2$ and $v_{0i}=0$. Reconstructing the corresponding
$\sigma$ by \eqref{geometry:sigma} in the original tensor metric
equation~\eqref{geometry:metriccomponents} yields
\begin{equation}\label{aw:residual}
 \frac{E_{0ii}}{h_i}
 =(Q_i-1)+2v_{i0}-v_{i0}^{\,2}
 =-\frac{(Q_i-1)^2}{4}\ne0.
\end{equation}
Here $Q_i-1$ is a unit on our metric-ratio base. Thus the unique
connection for the Aschieri--Weber predicate does not satisfy our
original full tensor metric predicate under this same-metric
left/right identification. This establishes the stated comparison;
it excludes no arbitrary alternative realization of the geometric data.

The standard-library checker
\path{literature/comparison_checks/aw_comparison_check.py}
recomputes both inverse compositions at $q^0$, the rational-tensor
identity on every vertex basis function and direction, and the six
normalized residuals on the two slices. Its optional
\texttt{--source-pdf} argument identifies the inspected arXiv v2 bytes;
it is not needed for the arithmetic. This bounded comparison check
is separate from the complete $512/204$ counting replay. The general
identities and their geometric interpretation are the written
derivation above, not an inference from a software status.

% END SHARED AW COMPARISON

\paragraph{Data availability.}
The exact coefficient data, verification programs, certificates and
execution records are publicly available in the versioned GitHub research
package~\cite{qlc-data}, at commit \texttt{1cefff4}.
\href{https://github.com/Torsheim/quantum-gravity-qlc/blob/1cefff4a809dae7626169fad6f469b97babf8a95/Quantum_Gravity_QLC_Computational_Supplement.zip}{The complete code/data archive}
accompanies article revision 0.13 and supplement revision 2.3.
It contains the same $175$ protected counting inputs as the delivered
supplement. The current package also includes the complete editable
sources, pinned requirements and exact replay records.
Run \path{reproduce.py}; \path{inputs.json} and
\path{synchronization.json} identify and bind the proof inputs and sources.

\paragraph{Funding.}
This research received no external funding.

\paragraph{Declaration of competing interest.}
The author declares no known competing financial interests or personal
relationships that could have appeared to influence the work reported in
this paper.

\section*{Declaration of generative AI and AI-assisted technologies
in the manuscript preparation process}
During the preparation of this manuscript and its computational supplement,
the author used ChatGPT and Codex (OpenAI) for drafting and revising the
text, mathematical support, and internal AI-assisted review of the arguments
and computations. Further details are given in the reproducibility section.
The author previously reviewed the manuscript and supplement, authorized
the subsequent documentation revisions, and takes full responsibility for
the content.

\begin{thebibliography}{99}\small
\bibitem{beggs-majid-star} E. J. Beggs and S. Majid,
\emph{*-Compatible Connections in Noncommutative Riemannian Geometry},
Journal of Geometry and Physics \textbf{61} (2011), 95--124.
Cited version: \url{https://arxiv.org/abs/0904.0539v2}.
\bibitem{majid-graphs} S. Majid,
\emph{Noncommutative Riemannian geometry on graphs},
Journal of Geometry and Physics \textbf{69} (2013), 74--93.
Cited version: \url{https://arxiv.org/abs/1011.5898v4}.
\bibitem{stacks-dvr} The Stacks Project, Lemma 32.15.1, Tag 0CM2,
a discrete-valuation trait reaching a point in a finite-type closure.
\url{https://stacks.math.columbia.edu/tag/0CM2}.
\bibitem{stacks-qf} The Stacks Project, Lemma 29.57.2, Tag 01TI,
openness of the quasi-finite locus.
\url{https://stacks.math.columbia.edu/tag/01TI}.
\bibitem{stacks-square} The Stacks Project, Lemma 49.10.1, Tag 0BWE,
quasi-finite localized square presentations and flatness.
\url{https://stacks.math.columbia.edu/tag/0BWE}.
\bibitem{stacks-jacobian} The Stacks Project, Lemma 29.37.15, Tag 02GU,
the square Jacobian criterion for \'etaleness.
\url{https://stacks.math.columbia.edu/tag/02GU}.
\bibitem{stacks-proper-qf} The Stacks Project, Lemma 37.44.1, Tag 02LS,
proper locally quasi-finite morphisms are finite.
\url{https://stacks.math.columbia.edu/tag/02LS}.
\bibitem{stacks-finite-free} The Stacks Project, Section 29.49, Tag 02K9,
finite locally free morphisms and locally constant rank.
\url{https://stacks.math.columbia.edu/tag/02K9}.
\bibitem{stacks-closure} The Stacks Project, Section 29.6, Tag 01R5,
scheme-theoretic image and restriction to opens.
\url{https://stacks.math.columbia.edu/tag/01R5}.
\bibitem{orig-emergence} S. Majid,
\emph{On the emergence of the structure of Physics},
preprint, cited version: \url{https://arxiv.org/abs/1711.00556v2}.
\bibitem{orig-square} S. Majid,
\emph{Quantum gravity on a square graph},
Classical and Quantum Gravity \textbf{36} (2019), 245009,
doi:10.1088/1361-6382/ab4975.
Cited version: 2018 preprint, \url{https://arxiv.org/abs/1810.10831v1}.
\bibitem{orig-book} E. J. Beggs and S. Majid,
\emph{Quantum Riemannian Geometry}, Grundlehren der mathematischen
Wissenschaften \textbf{355}, Springer, 2020.
\url{https://doi.org/10.1007/978-3-030-30294-8}.
Inspected: Example~8.20, pp.~594--598, with its relevant calculus,
connection, tensor, star and group-component dependencies; especially
Definition~8.3, Proposition~8.11 and Lemma~8.17.
\bibitem{orig-digital} S. Majid and A. Pacho\l{},
\emph{Digital finite quantum Riemannian geometries},
Journal of Physics A: Mathematical and Theoretical \textbf{53} (2020), 115202.
Cited version: \url{https://arxiv.org/abs/1807.08492v2}.
\bibitem{orig-polygon} J. N. Argota-Quiroz and S. Majid,
\emph{Quantum gravity on polygons and $\mathbb R\times\mathbb Z_n$ FLRW model},
Classical and Quantum Gravity \textbf{37} (2020), 245001.
Cited version: \url{https://arxiv.org/abs/2005.13999v3}.
\bibitem{orig-mr25} B. Mesland and A. Rennie,
\emph{Existence and uniqueness of the Levi-Civita connection on
noncommutative differential forms},
Advances in Mathematics \textbf{468} (2025), 110207.
Cited version: \url{https://arxiv.org/abs/2403.13735v3}.
\bibitem{orig-vh25} K. van Helden,
\emph{New Geometry in Classical and Quantum Gravity}, doctoral thesis,
University of Groningen, 2025, doi:10.33612/diss.1275744976.
\url{https://pure.rug.nl/ws/portalfiles/portal/1275744978/Complete_thesis.pdf}.
\bibitem{orig-graphs24} E. J. Beggs and S. Majid,
\emph{Quantum geodesic flows on graphs},
Letters in Mathematical Physics \textbf{114} (2024), 112.
\url{https://doi.org/10.1007/s11005-024-01860-6}.
\bibitem{orig-bgg} J. Bhowmick, B. Ghosh and S. Guin,
\emph{On two notions of torsion and metric compatibility of connections
in noncommutative geometry}, preprint (2025).
Cited version: \url{https://arxiv.org/abs/2509.11888v1}.
\bibitem{orig-fmr26} A. Flamant, B. Mesland and A. Rennie,
\emph{Comparison of Levi-Civita connections in noncommutative geometry},
Journal of Geometry and Physics \textbf{219} (2026), 105690,
doi:10.1016/j.geomphys.2025.105690.
Cited version: \url{https://arxiv.org/abs/2505.18692v1}.
\bibitem{orig-ah26} J. Arnlind and V. Hildebrandsson,
\emph{On the existence of noncommutative Levi-Civita connections in
derivation based calculi},
Journal of Geometry and Physics \textbf{222} (2026), 105764,
doi:10.1016/j.geomphys.2026.105764.
Cited version: \url{https://arxiv.org/abs/2505.13984v1}.
\bibitem{orig-aschieri26} P. Aschieri,
\emph{Cartan structure equations and Levi-Civita connection in
noncommutative geometry},
Journal of Geometry and Physics \textbf{224} (2026), 105800,
doi:10.1016/j.geomphys.2026.105800.
Cited version: \url{https://arxiv.org/abs/2006.02761v3}.
\bibitem{orig-aw25} P. Aschieri and T. Weber,
\emph{Metric compatibility and Levi-Civita Connections on Quantum Groups},
Journal of Algebra \textbf{661} (2025), 479--544,
doi:10.1016/j.jalgebra.2024.07.039.
Cited version: \url{https://arxiv.org/abs/2209.05453v2}.
\bibitem{qlc-data} [dataset] M. Torsheim,
\emph{Exact code and data for a marked path--triangle quantum
Levi--Civita classification}, GitHub, article revision 0.13 and
computational supplement revision 2.3 (2026).
Version: commit \texttt{1cefff4}.
\url{https://github.com/Torsheim/quantum-gravity-qlc/tree/1cefff4a809dae7626169fad6f469b97babf8a95}
(accessed 31 October 2026).
\end{thebibliography}
\end{document}
